\subsection{Integrals of operators}

Let G and H be two Hausdorff locally convex spaces over R, and suppose now that F is the space \(\mathcal{L}(\mathrm{G};\mathrm{H})\) of continuous linear mappings of G into H, equipped with the topology of pointwise convergence. Every continuous linear form on F may be extended to a continuous linear form on the product space \(H^{G}\) (TVS, II, §4, No.~1, Prop. 2), hence may be written \(u \mapsto \sum_{i=1}^{n} \langle u(\mathbf{a}_{i}), \mathbf{b}_{i}^{\prime} \rangle\) , where the \(a_{i}\) (resp. the \(b_{i}^{\prime}\) ) are elements of G (resp. of the dual \(H^{\prime}\) of H). To say that a mapping U of T into F is scalarly essentially \(\mu\) -integrable means that, for every \(a \in G\) and every \(b \in H^{\prime}\) , the numerical function \(t \mapsto \langle U(t) \cdot a, b^{\prime} \rangle\) is essentially \(\mu\) -integrable.

PROPOSITION 9. --- Let U be a scalarly essentially \(\mu\) -integrable mapping of T into \(\mathbf{F} = \mathcal{L}_{s}(\mathbf{G}; \mathbf{H})\) . In order that \(\int U d\mu \in F\) , it is necessary and sufficient that the following two conditions be satisfied:

\begin{enumerate}
\def\labelenumi{\alph{enumi})}
\item
  For every \(\mathbf{x} \in \mathbf{G}\) , \(\int (U(t) \cdot \mathbf{x}) d\mu(t) \in \mathrm{H}\) .
\item
  For every equicontinuous subset \(B'\) of \(H'\) , the set of linear forms \(u_{y'} : x \mapsto \int \langle U(t) \cdot x, y' \rangle d\mu(t)\) , where \(y'\) runs over \(B'\) , is equicontinuous.
\end{enumerate}

The conditions a) and b) are necessary. For, since for every \(\mathbf{x} \in G\) , the mapping \(\widetilde{\mathbf{x}}: V \mapsto V \cdot \mathbf{x}\) of \(\mathcal{L}_s(\mathrm{G};\mathrm{H})\) into H is linear, one sees (No.~1,

Prop. 1) that \(\widetilde{\mathbf{x}}\circ U:t\mapsto U(t)\cdot \mathbf{x}\) is scalarly essentially \(\mu\) -integrable and that

\[
S \cdot \mathbf {x} = \int (U (t) \cdot \mathbf {x}) d \mu (t),\tag{1}
\]

where \(S = \int U d\mu \in \mathcal{L}_s(\mathrm{G};\mathrm{H})\) . This proves a). Moreover, (1) may also be written

\[
\langle S \cdot \mathbf {x}, \mathbf {y} ^ {\prime} \rangle = \int \left\langle U (t) \cdot \mathbf {x}, \mathbf {y} ^ {\prime} \right\rangle d \mu (t) = \left\langle \mathbf {x}, u _ {\mathbf {y} ^ {\prime}} \right\rangle ,\tag{2}
\]

in other words \(^{t}S \cdot y' = u_{y'}\) . Since S is continuous, \(^{t}S\) transforms every equicontinuous subset of \(H'\) into an equicontinuous subset of \(G'\) , whence b).

Conversely, suppose a) and b) verified. By virtue of a), the formula (1) defines a linear mapping S of G into H, and, for every \(y' \in H'\) , this mapping satisfies (2) (No.~1, Prop. 1); but then, condition (b) says that S is continuous (TVS, II, §6, No.~4, Props. 5 and 6, and III, §3, No.~5, Prop. 7), therefore \(S \in \mathcal{L}_{s}(G;H)\) . Finally, formula (2) proves that \(S = \int U d\mu\) .

COROLLARY. --- The condition b) of Prop. 9 is satisfied in each of the following two cases:

\(1^{\circ}\) The measure \(\mu\) is bounded, and if S is its support, then \(U(S)\) is an equicontinuous subset of \(\mathcal{L}(G;H)\) .

\(2^{\circ}\) The condition a) of Prop. 9 is satisfied, the space G is barreled, and, for every compact subset K of T, \(U(\mathrm{K})\) is a bounded subset of \(\mathcal{L}_s(\mathrm{G};\mathrm{H})\) .

Let us first place ourselves in case \(1^{\circ}\) . We can restrict ourselves to the case that S = T (Ch. V, §7, No.~1, Th. 1). Then, for every equicontinuous subset \(B'\) of \(H'\) , there exists an equicontinuous, convex, balanced and weakly closed subset \(A' \subset G'\) such that \({}^{t}U(t) \cdot y' \in A'\) for all \(y' \in B'\) and all \(t \in T\) (TVS, II, §6, No.~4, Prop. 6). Since U is scalarly \(\mu\) -integrable, the mapping \(t \mapsto {}^{t}U(t) \cdot y'\) of T into the dual \(G'\) of G equipped with \(\sigma(G', G)\) , is scalarly \(\mu\) -integrable, and one may write

\[
u _ {\mathbf {y} ^ {\prime}} = \int \left(^ {t} U (t) \cdot \mathbf {y} ^ {\prime}\right) d \mu (t).
\]

Since \(A'\) is convex and compact for \(\sigma(G', G)\) , the Cor. of Prop. 5 of No.~2 shows that \(u_{\mathbf{y}'} \in \mu(T)A'\) for every \(\mathbf{y}' \in B'\) , which proves our assertion.

Let us now place ourselves in case \(2^{\circ}\) . For every \(y' \in H'\) and every compact subset K of T, set

\[
u _ {\mathrm{K}, \mathbf {y} ^ {\prime}} = \int \varphi_ {\mathrm{K}} (t) \left(^ {t} U (t) \cdot \mathbf {y} ^ {\prime}\right) d \mu (t),
\]

an element of the algebraic dual \(\mathbf{G}^*\) of \(\mathbf{G}\) . Since \(\mathbf{G}\) is barreled, every bounded subset of \(\mathcal{L}_s(\mathbf{G};\mathbf{H})\) is equicontinuous (TVS, III, §4, No.~2, Th. 1); the first part of the argument, applied to the function \(\varphi_{\mathbf{K}}U\) and the bounded measure \(\varphi_{\mathbf{K}}\cdot \mu\) , shows that \(u_{\mathbf{K},\mathbf{y}'}\in \mathbf{G}'\) . Moreover, for the topology \(\sigma (\mathbf{G}^{*},\mathbf{G})\) , one has \(u_{\mathbf{y}'} = \lim_{\mathbf{K}} = u_{\mathbf{K},\mathbf{y}'}\) , the limit being taken with respect to the increasing directed set of compact subsets of \(\mathbf{T}\) (Ch. V, §1, No.~3, Prop. 10). To verify condition b) of Prop. 9, it suffices, by Prop. 9, to prove that the linear mapping \(S\) of \(\mathbf{G}\) into \(\mathbf{H}\) defined by (1) is continuous; moreover, \(\mathbf{G}\) is barreled, therefore it will suffice to prove that \(S\) is continuous when \(\mathbf{G}\) and \(\mathbf{H}\) are equipped with their weakened topologies (TVS, IV, §1, No.~3, Prop. 7); finally, by virtue of (2), one is reduced to showing that \(u_{\mathbf{y}'}\in \mathbf{G}'\) for every \(\mathbf{y}'\in \mathbf{H}'\) . Since \(u_{\mathbf{y}'}\) is in the closure, for \(\sigma (\mathbf{G}^{*},\mathbf{G})\) , of the set \(\mathbf{M}'\) of the \(u_{\mathbf{K},\mathbf{y}'}\) , where \(\mathbf{K}\) runs over the set of compact subsets of \(\mathbf{T}\) , it suffices to prove that \(\mathbf{M}'\) is equicontinuous (TVS, III, §3, No.~4, Prop. 4); and since \(\mathbf{G}\) is barreled, it comes to the same to say that for every \(\mathbf{x}\in \mathbf{G}\) , the set of \(\langle \mathbf{x},u_{\mathbf{K},\mathbf{y}'}\rangle\) is bounded (TVS, III, §4, No.~2, Th. 1). But this follows at once from the relations

\[
| \langle \mathbf {x}, u _ {\mathrm{K}, \mathbf {y} ^ {\prime}} \rangle | = \left| \int \varphi_ {\mathrm{K}} (t) \langle U (t) \cdot \mathbf {x}, \mathbf {y} ^ {\prime} \rangle d \mu (t) \right| \leqslant \int | \langle U (t) \cdot \mathbf {x}, \mathbf {y} ^ {\prime} \rangle | d \mu (t).
\]

PROPOSITION 10. --- Let U be a mapping of T into F = \(\mathcal{L}_s(\mathrm{G};\mathrm{H})\) . In each of the following three cases, U is scalarly essentially \(\mu\) -integrable and \(\int U d\mu \in \mathcal{L}_s(\mathrm{G};\mathrm{H})\) :

\begin{enumerate}
\def\labelenumi{\alph{enumi})}
\item
  H is quasi-complete, \(\mu\) is bounded and, if S is its support, U is \(\mu\) -measurable and \(U(S)\) is equicontinuous.
\item
  H is semi-reflexive, \(\mu\) is bounded and, if S is its support, U is scalarly \(\mu\) -measurable and \(U(S)\) is equicontinuous.
\item
  H is semi-reflexive, G is barreled, U is scalarly essentially \(\mu\) -integrable, and, for every compact subset K of T, \(U(\mathrm{K})\) is bounded.
\end{enumerate}

The fact that U is scalarly essentially integrable is obvious in all three cases; by virtue of Prop. 9 and its corollary, it suffices in each case to verify condition a) of Prop. 9. Now, this condition follows from Prop. 8 of No.~2 in the first case, and from the Cor. of Prop. 7 of No.~2 in the other two cases.

\subsection{The property (GDF)}

In this No.~we are going to consider locally convex spaces F having the following property (called the `countably closed graph' property): \(^{6}\)

(GDF) If u is a linear mapping of F into a Banach space B such that, in the product space \(F \times B\) , every limit of every convergent sequence of points of the graph \(\Gamma\) of u again belongs to \(\Gamma\) , then u is continuous.

Every Fréchet space has the property (GDF) (TVS, I, §3, No.~3, Cor. 5 of Th. 1). In the Appendix, we shall see other examples of spaces having the property (GDF).

PROPOSITION 11. --- Every Hausdorff locally convex space F with the property (GDF) is barreled.

Let V be a barrel in F, q is gauge, which is a semi-norm on F; let H be the Hausdorff space associated with the space F equipped with the topology defined by this single semi-norm. The completion \(\hat{H}\) of H is a Banach space; let \(\pi\) be the canonical mapping of F into \(\hat{H}\) ; we are going to show that \(\pi\) is continuous (for the original topology on F); this will establish the proposition, because V, the inverse image under \(\pi\) of the unit ball of \(\hat{H}\) , will then be a neighborhood of 0 in F. To establish the continuity of \(\pi\) it will suffice, by virtue of (GDF), to show that the graph of \(\pi\) is closed in \(F \times \hat{H}\) ; in other words, we must see that if \(\mathfrak{F}\) is a filter on F, convergent to \(x \in F\) , and its image \(\pi(\mathfrak{F})\) converges to \(y \in \hat{H}\) , then \(y = \pi(x)\) . Now, every element \(x'\) of the polar \(V^{\circ}\) of V in \(F'\) may be extended in only one way to a continuous linear form on \(\hat{H}\) (again denoted \(x'\) ), and the set of these forms is the unit ball of the dual of \(\hat{H}\) ; it therefore suffices to show that \(\langle y, x' \rangle = \langle \pi(x), x' \rangle\) for every \(x' \in V^{\circ}\) . But this follows from the relations

\[
\langle \mathbf {y}, \mathbf {x} ^ {\prime} \rangle = \lim _ {\mathfrak {F}} \left\langle \pi (\mathbf {z}), \mathbf {x} ^ {\prime} \right\rangle = \lim _ {\mathfrak {F}} \left\langle \mathbf {z}, \mathbf {x} ^ {\prime} \right\rangle = \left\langle \mathbf {x}, \mathbf {x} ^ {\prime} \right\rangle = \left\langle \pi (\mathbf {x}), \mathbf {x} ^ {\prime} \right\rangle .
\]

THEOREM 1. (Gelfand--Dunford)--- Let F be a Hausdorff locally convex space having the property (GDF), \(F_{s}^{\prime}\) its weak dual. For every mapping f of T into \(F_{s}^{\prime}\) that is scalarly essentially \(\mu\) -integrable, the integral \(\int f d\mu\) belongs to \(F^{\prime}\) .

Recall that the dual of \(F_{s}^{\prime}\) is F (TVS, II, §6, No.~2, Prop. 3). For every \(z \in F\) , the numerical function \(\langle z, f \rangle\) is therefore essentially \(\mu\) -integrable; let \(\theta(z)\) be its class in \(L^{1}(\mu)\) . To show that \(\int f d\mu \in F^{\prime}\) , one must establish that the linear form \(z \mapsto \langle z, \int f d\mu \rangle\) is continuous on F; in fact, we are going to prove the following stronger result:

Lemma 2. --- Let f be a mapping of T into \(F_{s}^{\prime}\) , such that for every \(z \in F\) , the numerical function \(\langle z, f \rangle\) belongs to \(\overline{\mathcal{L}^{p}}(\mu) (1 \leqslant p \leqslant +\infty)\) ; let \(\theta(z)\) be the class of this function in \(L^{p}(\mu)\) . Then \(z \mapsto \theta(z)\) is a continuous linear mapping of F into \(L^{p}(\mu)\) .

By virtue of the (GDF) property, it suffices to show that for every sequence \((\mathbf{z}_{n})\) of elements of F converging to z, such that \((\theta(\mathbf{z}_{n}))\) converges to \(u \in \mathrm{L}^{p}(\mu)\) , one has \(u = \theta(\mathbf{z})\) . Now, replacing if necessary the sequence \((\mathbf{z}_{n})\)

by a subsequence, one can suppose that the sequence of functions \(\langle \mathbf{z}_n,\mathbf{f}\rangle\) converges locally almost everywhere to a function \(h\in \overline{\mathcal{L}^p} (\mu)\) , with class \(u\) in \(\mathrm{L}^p (\mu)\) (Ch. IV, §3, No.~4, Th. 3 and Ch. V, §1, No.~3). Since, by hypothesis, for every \(t\in \mathrm{T}\) the sequence \((\langle \mathbf{z}_n,\mathbf{f}(t)\rangle)\) converges to \(\langle \mathbf{z},\mathbf{f}(t)\rangle\) , we have \(h(t) = \langle \mathbf{z},\mathbf{f}(t)\rangle\) locally almost everywhere, consequently \(u = \theta (\mathbf{z})\) .

COROLLARY 1. --- Let \(G_{i}\) ( \(1 \leqslant i \leqslant n\) ) be n Hausdorff locally convex spaces having the property (GDF), and let F be the space of separately continuous multilinear forms on \(\prod_{i=1}^{n} G_{i}\) , equipped with the topology of pointwise convergence. For every mapping f of T into F that is scalarly essentially \(\mu\) -integrable, one has \(\int f d\mu \in F\) .

The space F is in duality with the tensor product \(\bigotimes_{i=1}^{n}G_{i}\) , and the topology of pointwise convergence on F is none other than the topology \(\sigma(F,\bigotimes_{i=1}^{n}G_{i})\) . The algebraic dual \(F^{\prime*}\) is therefore the space of all multilinear forms on \(\prod_{i=1}^{n}G_{i}\) . Let \(\mathbf{z}=(\mathbf{z}_{1},\ldots,\mathbf{z}_{n})\) be an element of \(\prod_{i=1}^{n}G_{i}\) ; for every multilinear form \(u\in F^{\prime*}\) , the mapping \(\mathbf{x}\mapsto\mathbf{u}(\mathbf{z}_{1},\ldots,\mathbf{z}_{i-1},\mathbf{x},\mathbf{z}_{i+1},\ldots,\mathbf{z}_{n})\) is a linear form on \(G_{i}\) , which we shall denote by \(\lambda_{i}(\mathbf{z})(\mathbf{u})\) ; we thus obtain a linear mapping \(\lambda_{i}(\mathbf{z})\) of \(F^{\prime*}\) into the algebraic dual \(G_{i}^{*}\) of \(G_{i}\) , continuous for the topologies \(\sigma(F^{\prime*},\bigotimes_{i=1}^{n}G_{i})\) and \(\sigma(G_{i}^{*},G_{i})\) . To say that \(u\in F\) means that for every index i and every \(z\in\Pi_{i=1}^{n}G_{i}\) , one has \(\lambda_{i}(\mathbf{z})(\mathbf{u})\in G_{i}^{\prime}\) . Now, by Prop. 1 of No.~1, the mapping \(\lambda_{i}(\mathbf{z})\circ\mathbf{f}\) is a scalarly essentially \(\mu\) -integrable mapping of T into \(G_{i}^{\prime}\) equipped with the topology \(\sigma(G_{i}^{\prime},G_{i})\) , and

\[
\int \left(\lambda_ {i} (\mathbf {z}) \circ \mathbf {f}\right) d \mu = \lambda_ {i} (\mathbf {z}) \left(\int \mathbf {f} d \mu\right).
\]

By Th. 1, \(\int (\lambda_i(\mathbf{z})\circ \mathbf{f})d\mu \in \mathrm{G}_i^\prime\) for \(1\leqslant i\leqslant n\) , therefore \(\int \mathbf{f}d\mu \in \mathbf{F}\)

COROLLARY 2. --- Let G be a Hausdorff locally convex space having the property (GDF), and H a semi-reflexive space whose strong dual \(H_{b}^{\prime}\) has the property (GDF) (cf.~App., No.~2, Prop. 3). Let F be the space \(\mathcal{L}_{s}(G;H)\) ; for every mapping U of T into F that is scalarly essentially \(\mu\) -integrable, the integral \(\int U d\mu\) belongs to F.

Since G is barreled (Prop. 11), \(\mathcal{L}(\mathrm{G};\mathrm{H}) = \mathcal{L}(\mathrm{G}_{\sigma};\mathrm{H}_{\sigma})\) (TVS, IV, §1, No.~3, Prop. 7); moreover, one can replace \(\mathrm{F} = \mathcal{L}_s(\mathrm{G};\mathrm{H})\) by the space \(\mathcal{L}_s(\mathrm{G}_{\sigma};\mathrm{H}_{\sigma})\) , since the two spaces have the same dual \(\mathrm{G} \otimes \mathrm{H}'\) (TVS, II, §6, No.~2, Prop. 3, and the first paragraph of No.~3 above). If, for every \(u \in \mathcal{L}(\mathrm{G};\mathrm{H}) = \mathcal{L}(\mathrm{G}_{\sigma};\mathrm{H}_{\sigma})\) , one sets \(\widetilde{u}(\mathbf{x},\mathbf{y}') = \langle u(\mathbf{x},\mathbf{y}') \rangle\) (for \(\mathbf{x} \in \mathbf{G}\) , \(y' \in H')\) , the linear mapping \(u \mapsto \widetilde{u}\) is a bijection of F onto the space \(F_1\) of separately continuous bilinear forms on \(G_\sigma \times H_s'\) , where \(H_s'\) denotes the dual \(H'\) equipped with the weak topology \(\sigma(H', H)\) (App., No.~1); moreover, this mapping is an isomorphism of \(\mathcal{L}_s(G_\sigma; H_\sigma)\) onto \(F_1\) equipped with the topology of pointwise convergence (loc. cit.). But since by hypothesis H is the dual of \(H_b'\) , \(F_1\) is also the space of separately continuous bilinear forms on \(G \times H_b'\) . Cor. 2 therefore follows from Cor. 1.

Note that Cor. 2 is applicable in particular when G is Banach space and H is a reflexive Banach space.

\subsection{Measurable mappings and scalarly measurable mappings}

If a mapping f of T into a Hausdorff locally convex space F is scalarly \(\mu\) -measurable, it does not in general follow that f is \(\mu\) -measurable (Exer. 12). Nevertheless:

PROPOSITION 12. --- If F is a separable, metrizable locally convex space, then every scalarly \(\mu\) -measurable mapping f of T into F is \(\mu\) -measurable.

For, F may be regarded as a subspace of a countable product \(\prod_{n}E_{n}\) of Banach spaces (TVS, II, §1, No.~3, Prop. 3), and one can suppose that \(\mathrm{pr}_n(\mathbf{F})\) is dense in \(\mathbf{E}_n\) , which is therefore separable. For every \(n\) , the mapping \(\mathrm{pr}_n \circ \mathbf{f}\) is scalarly \(\mu\) -measurable, hence \(\mu\) -measurable (Ch. IV, §5, No.~5, Cor. 2 of Prop. 10), consequently \(\mathbf{f}\) is \(\mu\) -measurable (Ch. IV, §5, No.~3, Th. 1).

PROPOSITION 13. --- Let F be a locally convex space that is the direct limit of a sequence of separable, metrizable locally convex spaces \(F_{n}\) , F being the union of the \(F_{n}\) . Let \(F'\) be the dual of F, equipped with the topology \(\sigma(\mathbf{F}',\mathbf{F})\) . Then, every scalarly \(\mu\) -measurable mapping f of T into \(F'\) is \(\mu\) -measurable.

Suppose first that F is metrizable and separable, and let D be a countable dense set in F. Let \((\mathrm{V}_{n})\) be a decreasing fundamental sequence of balanced, convex, open neighborhoods of 0 in F; the polar sets \(V_{n}^{\circ}\) are equicontinuous and their union is all of \(F'\) . Let \(\mathrm{T}_{n} = \mathbf{f}^{-1}(\mathrm{V}_{n}^{\circ})\) ; the sequence \((\mathrm{T}_{n})\) is increasing and \(T = \bigcup_{n} T_{n}\) ; let us show that each of the \(T_{n}\) is \(\mu\) -measurable. Indeed, \(D \cap V_{n}\) is dense in \(V_{n}\) ; for every \(y \in D \cap V_{n}\) , let \(S_{y}\) be the set of \(t \in T\) such that \(|\langle y, f(t) \rangle| \leqslant 1\) ; the hypothesis implies that each of the \(S_{y}\) is measurable, and \(T_{n}\) is the intersection of the countable family of the \(S_{y}\) \((y \in D \cap V_{n})\) . This being so, for every compact subset K of T and every \(\varepsilon > 0\) , there exists an integer n such that \(\mu\left(\mathrm{K}-\left(\mathrm{K}\cap\mathrm{T}_{n}\right)\right)\leqslant\frac{\varepsilon}{4}\) , then a compact subset \(K_{1}\) of \(K\cap T_{n}\) such that \(\mu\left(\left(\mathrm{K}\cap\mathrm{T}_{n}\right)-\mathrm{K}_{1}\right)\leqslant\frac{\varepsilon}{4}\) ; finally, there exists a compact subset \(K_{2}\) of \(K_{1}\) such that \(\mu(K_{1}-K_{2})\leqslant\frac{\varepsilon}{2}\) and such that the restrictions to \(K_{2}\) of all of the functions \(\langle y,f\rangle\) , where \(y\in D\) , are continuous (Ch. IV, §5, No.~1, Prop. 2). Since the set \(\mathbf{f}(K_{2})\subset\mathbf{f}(T_{n})\subset V_{n}^{\circ}\) is equicontinuous, the topology induced by \(\sigma(F',F)\) on \(\mathbf{f}(K_{2})\) is identical to the topology of pointwise convergence in D (GT, X, §2, No.~4, Th. 1); consequently, the restriction of f to \(K_{2}\) is continuous, whence our assertion in the first case.

Let us pass to the general case. If \(z^{\prime}\) is a continuous linear form on F, its restriction \(z_{n}^{\prime}\) to \(F_{n}\) is continuous; since \(F = \bigcup_{n} F_{n}\) , the dual \(F^{\prime}\) of F may be identified (algebraically) with a linear subspace of the product \(\prod_{n} F_{n}^{\prime}\) , and \(pr_{n}z^{\prime} = z_{n}^{\prime}\) . Moreover, since each finite subset of F is contained in one of the \(F_{n}\) , the topology \(\sigma(F^{\prime}, F)\) is none other than the topology induced by the product topology of the topologies \(\sigma(F_{n}^{\prime}, F_{n})\) . This being so, if f is scalarly \(\mu\) -measurable then \(pr_{n} \circ f\) is scalarly \(\mu\) -measurable, since for every \(t \in T\) , \(\operatorname{pr}_{n}(f(t))\) is the restriction of \(f(t)\) to \(F_{n}\) . The first part of the proof shows that \(pr_{n} \circ f\) is \(\mu\) -measurable for every n, therefore so is f (Ch. IV, §5, No.~3, Th. 1).

\subsection{Applications: I. Extension of a continuous function to a space of measures}

Let T be a locally compact space, F a quasi-complete, Hausdorff locally convex space, and f a continuous mapping of T into F; if \(\mu\) is a positive measure on T, with compact support S, then \(\mathbf{f}(S)\) is compact; the closed convex envelope of \(\mathbf{f}(S)\) is then compact (TVS, III, §1, No.~6), therefore f is scalarly \(\mu\) -integrable and \(\int f d\mu \in F\) (No.~2, Cor. of Prop. 5). If now \(\lambda\) is any real measure with compact support, \(\lambda^{+}\) and \(\lambda^{-}\) are positive measures with compact support; setting \(\int f d\lambda = \int f d\lambda^{+} - \int f d\lambda^{-}\) , one verifies immediately (using the relation \((\lambda + \mu)^{+} + \lambda^{-} + \mu^{-} = \lambda^{+} + \mu^{+} + (\lambda + \mu)^{-}\) ) that \(\lambda \mapsto \int f d\lambda\) is a linear mapping of the space \(\mathcal{C}'(T)\) of measures on T with compact support, into the locally convex space F.

Let us now observe that the space \(\mathcal{C}^{\prime}(\mathrm{T})\) may be identified with the dual of the space \(\mathcal{C}(\mathrm{T})\) of continuous numerical functions on T (whence its notation), when \(\mathcal{C}(\mathrm{T})\) is equipped with the topology of compact convergence (which we shall always assume in this No.~and in the following one): for, it is known on the one hand (Ch. IV, §4, No.~8, Prop. 14) that the measures on T that can be extended to continuous linear forms on \(\mathcal{C}(\mathrm{T})\) are the measures with compact support, and conversely, the restriction to \(\mathcal{K}(T)\) of a continuous linear form on \(\mathcal{C}(T)\) is a measure (since the topology of \(\mathcal{K}(T)\) is finer than the one induced by the topology of \(\mathcal{C}(T)\) ).

PROPOSITION 14. --- Let T be a locally compact space, F a quasi-complete, Hausdorff locally convex space, and f a continuous mapping of T into F. If the space \(\mathcal{C}'(T)\) of measures on T with compact support is equipped with the topology of uniform convergence in the compact subsets of \(\mathcal{C}(T)\) , then the mapping \(\lambda \mapsto \int \mathbf{f} d\lambda\) is the unique continuous linear mapping \(\widetilde{\mathbf{f}}\) of \(\mathcal{C}'(T)\) into F such that \(\widetilde{\mathbf{f}}(\varepsilon_t) = \mathbf{f}(t)\) for every \(t \in T\) .

To establish the uniqueness of the extension, it suffices to see that the point measures \(\varepsilon_{t}\) form a total set in \(\mathcal{C}'(T)\) ; since the dual of \(\mathcal{C}'(T)\) is \(\mathcal{C}(T)\) (TVS, IV, §1, No.~1, Th. 1), it suffices to observe that every function \(g \in \mathcal{C}(T)\) that is orthogonal to all of the measures \(\varepsilon_{t}\) is equal to 0 by definition (TVS, IV, §1, No.~2, Prop. 2).

Let us now show that \(\lambda \mapsto \int \mathbf{f}d\lambda\) is continuous. Let \(V\) be a closed, balanced convex neighborhood of 0 in \(F\) ; it suffices to prove that there exists a relatively compact subset \(L\) of \(\mathcal{C}(T)\) such that the relations \(\lambda \in L^{\circ}\) and \(\mathbf{z}' \in V^{\circ}\) imply \(\left|\left\langle \int \mathbf{f}d\lambda, \mathbf{z}' \right\rangle \right| \leqslant 1\) , or again that \(\left|\int \langle \mathbf{f}, \mathbf{z}' \rangle d\lambda \right| \leqslant 1\) . To this end, we are going to show that as \(\mathbf{z}'\) runs over \(V^{\circ}\) , the set \(L\) of numerical functions \(\langle \mathbf{f}, \mathbf{z}' \rangle\) is relatively compact in \(\mathcal{C}(T)\) . Since \(V^{\circ}\) is bounded for \(\sigma(F', F)\) , the supremum of the numbers \(|\langle \mathbf{f}(t), \mathbf{z}' \rangle|\) , for \(t \in T\) fixed and \(\mathbf{z}'\) running over \(V^{\circ}\) , is finite; by virtue of Ascoli's theorem (GT, X, §2, No.~5, Cor. 2 of Th. 2), it therefore suffices to show that the set of \(\langle \mathbf{f}, \mathbf{z}' \rangle\) ( \(\mathbf{z}' \in V^{\circ}\) ) is equicontinuous. But, for every \(t_0 \in T\) and every \(\delta > 0\) , there exists by hypothesis a neighborhood \(W\) of \(t_0\) in \(T\) such that \(\mathbf{f}(t) - \mathbf{f}(t_0) \in \delta V\) for all \(t \in W\) ; it follows that \(|\langle \mathbf{f}(t), \mathbf{z}' \rangle - \langle \mathbf{f}(t_0), \mathbf{z}' \rangle| \leqslant \delta\) for all \(t \in W\) and all \(\mathbf{z}' \in V^{\circ}\) , which completes the proof.

Remarks. --- 1) The mapping \(t \mapsto \varepsilon_t\) is a homeomorphism of T into the space \(\mathcal{C}'(\mathrm{T})\) ; for, if L is a compact subset of \(\mathcal{C}(\mathrm{T})\) , and \(t_0 \in \mathrm{T}\) , there exists (GT, X, §2, No.~5, Cor. 3 of Th. 2) a neighborhood W of \(t_0\) such that \(|g(t) - g(t_0)| \leqslant 1\) for every \(t \in \mathrm{W}\) and every function \(g \in \mathrm{L}\) , therefore \(\varepsilon_t - \varepsilon_{t_0} \in \mathrm{L}^\circ\) for \(t \in \mathrm{W}\) , which proves the continuity of \(t \mapsto \varepsilon_t\) (cf.~Ch. IV, §4, No.~8, Prop. 15); one knows in addition that the inverse mapping is already continuous for the vague topology (Ch. III, §1, No.~9, Prop. 13), hence a fortiori for the topology of uniform convergence in the compact subset of \(\mathcal{C}(\mathrm{T})\) . If one then identifies T with its image in \(\mathcal{C}'(\mathrm{T})\) via \(t \mapsto \varepsilon_t\) , one can say that \(\lambda \mapsto \int f d\lambda\) is the unique continuous extension of f to a linear mapping.

\begin{enumerate}
\def\labelenumi{\arabic{enumi})}
\setcounter{enumi}{1}
\tightlist
\item
  Note that in the proof of the continuity of \(\lambda\mapsto\int f d\lambda\) , we have not used the fact that F is quasi-complete. The conclusion of Prop. 14 is therefore still valid without this hypothesis, provided one knows in addition that \(\int f d\mu\in F\) for every positive measure \(\mu\) with compact support.
\end{enumerate}

Suppose now that \(\mathbf{f}(\mathrm{T})\) is a bounded subset of \(\mathrm{F}\) . Then, for every bounded positive measure \(\mu\) on \(\mathrm{T}\) , \(\mathbf{f}\) is scalarly \(\mu\) -integrable and \(\int \mathbf{f}d\mu \in \mathrm{F}\) (No.~2, Prop. 8). If \(\lambda\) is any bounded real measure on T, then \(\lambda^{+}\) and \(\lambda^{-}\) are bounded, and one sees immediately that \(\lambda \mapsto \int \mathbf{f} d\lambda\) , defined as above, is a linear mapping of the space \(\mathcal{M}^1(\mathrm{T})\) of bounded measures on T, into the locally convex space F, that obviously extends the mapping \(\lambda \mapsto \int \mathbf{f} d\lambda\) of \(\mathcal{C}'(\mathrm{T})\) into F.

PROPOSITION 15. --- Let T be a locally compact space, F a Hausdorff locally convex space that is quasi-complete, and f a continuous mapping of T into F such that \(\mathbf{f}(\mathrm{T})\) is bounded. If \(\mathcal{M}^{1}(\mathrm{T})\) is equipped with its Banach space topology, then the linear mapping \(\lambda \mapsto \int f d\lambda\) of \(\mathcal{M}^{1}(\mathrm{T})\) into F is continuous.

For every closed, balanced, convex neighborhood V of 0 in F, there exists a \(\rho > 0\) such that \(\mathbf{f}(\mathrm{T}) \subset \rho\mathrm{V}\) ; the closed, balanced, convex envelope B of \(\mathbf{f}(\mathrm{T})\) is therefore contained in \(\rho V\) , and it is complete by hypothesis. If then \(\|\lambda\| \leqslant 1/\rho\) , it follows from No.~2, Prop. 8, and the relation \(\|\lambda\| = \lambda^{+}(\mathrm{T}) + \lambda^{-}(\mathrm{T})\) , that \(\int f d\lambda \in B/\rho \subset V\) .

\subsection{Applications: II. Extension, to a space of measures, of a continuous function with values in a space of operators}

Let G be a Hausdorff locally convex space, H a Hausdorff and quasi-complete locally convex space, and denote by F the space \(\mathcal{L}(\mathrm{G};\mathrm{H})\) of continuous linear mappings of G into H, equipped with the topology of compact convergence. The space F is not necessarily quasi-complete, and if \(t \mapsto U(t)\) is a continuous mapping of T into F, and \(\mu\) is a positive measure on T with compact support, one does not necessarily have \(\int U d\mu \in F\) (Exer. 27). However, if, for every compact subset K of T, \(U(\mathrm{K})\) is equicontinuous, then its balanced convex envelope in F is also equicontinuous (TVS, III, §3, No.~4), and since H is quasi-complete, the closure of this balanced convex envelope will be a complete subset of F (TVS, III, §3, No.~8, Prop. 11); one will then indeed have \(\int U d\mu \in F\) (No.~2, Prop. 8).

The supplementary condition imposed on U may be expressed otherwise:

Lemma 3. --- Let G, H be two locally convex spaces, U a mapping of a locally compact space T into \(\mathcal{L}(\mathrm{G};\mathrm{H})\) . The following conditions are equivalent:

\begin{enumerate}
\def\labelenumi{\alph{enumi})}
\item
  The mapping \((t,\mathbf{x})\mapsto U(t)\cdot \mathbf{x}\) of \(\mathrm{T}\times \mathrm{G}\) into H is continuous.
\item
  For every compact subset K of T, \(U(\mathbf{K})\) is equicontinuous, and there exists a total set \(D \subset G\) such that for every \(x \in D\) , the mapping \(t \mapsto U(t) \cdot x\) is continuous on T.
\end{enumerate}

Moreover, when \(U\) satisfies these conditions, \(U\) is a continuous mapping of \(\mathrm{T}\) into \(\mathcal{L}(\mathrm{G};\mathrm{H})\) equipped with the topology of compact convergence.

To see that a) implies b), we observe that for every neighborhood V of 0 in H and every \(t \in \mathbf{K}\) , there exist by hypothesis a neighborhood \(\mathrm{L}_t\) of \(t\) in T and a neighborhood \(\mathrm{W}_t\) of 0 in G such that the relations \(t' \in \mathrm{L}_t\) and \(\mathbf{x} \in \mathrm{W}_t\) imply \(U(t') \cdot \mathbf{x} \in \mathrm{V}\) . It suffices to cover K with a finite number of neighborhoods \(\mathrm{L}_{t_i}\) and to take \(\mathrm{W} = \bigcap_{i} \mathrm{W}_{t_i}\) in order to have \(U(t) \cdot \mathbf{x} \in \mathrm{V}\) whenever \(t \in \mathrm{K}\) and \(\mathbf{x} \in \mathrm{W}\) , which proves the equicontinuity of \(U(\mathrm{K})\) .

Conversely, suppose b) is verified; it suffices to show that for every compact subset K of T, the mapping \((t,\mathbf{x})\mapsto U(t)\cdot\mathbf{x}\) is continuous on \(K\times G\) . Let \(M=U(K)\) ; since M is equicontinuous, it follows that on M, the topology of pointwise convergence in G is identical to the topology of pointwise convergence in D (GT, X, §2, No.~4, Th. 1); the hypothesis b) therefore implies that \(t\mapsto U(t)\) is a continuous mapping of K into \(\mathcal{L}(G;H)\) when \(\mathcal{L}(G;H)\) is equipped with the topology of pointwise convergence. On the other hand, \((A,\mathbf{x})\mapsto A\cdot\mathbf{x}\) is a continuous mapping of \(M\times G\) into H when M is equipped with the topology of pointwise convergence (GT, X, §2, No.~1, Cor. 4 of Prop. 1). Since the mapping \((t,\mathbf{x})\mapsto U(t)\cdot\mathbf{x}\) may be factored as \((t,\mathbf{x})\mapsto(U(t),\mathbf{x})\mapsto U(t)\cdot\mathbf{x}\) , we conclude that it is continuous.

Finally, the last assertion of the lemma follows from the fact that, on M, the topology of compact convergence is identical to that of pointwise convergence (GT, X, §2, No.~4, Th. 1).

Thus, suppose that U satisfies the conditions of Lemma 3; then (if H is quasi-complete) one defines, as in No.~6, a linear mapping \(\lambda \mapsto \int U d\lambda\) of \(\mathcal{C}'(\mathrm{T})\) into \(\mathrm{F} = \mathcal{L}(\mathrm{G}; \mathrm{H})\) . We set \(U(\lambda) = \int U d\lambda\) .

PROPOSITION 16. --- Let G, H be two Hausdorff locally convex spaces, with H assumed to be quasi-complete. Let U be a mapping of T into \(\mathcal{L}(\mathrm{G};\mathrm{H})\) such that \((t,\mathbf{x})\mapsto U(t)\cdot \mathbf{x}\) is a continuous mapping of \(\mathrm{T}\times \mathrm{G}\) into H. Then the bilinear mapping \((\lambda ,\mathbf{x})\mapsto U(\lambda)\cdot \mathbf{x}\) of \(\mathcal{C}'(\mathrm{T})\times \mathrm{G}\) into H is hypocontinuous relative to the equicontinuous subsets of \(\mathcal{C}'(\mathrm{T})\) and the compact subsets of G (which implies that the linear mapping \(\lambda \mapsto U(\lambda)\) of \(\mathcal{C}'(\mathrm{T})\) into F is continuous).

The continuity of \(\lambda \mapsto U(t)\) as a mapping of \(\mathcal{C}'(\mathrm{T})\) into \(\mathrm{F}\) follows from Lemma 3 and Remark 2 following Prop. 14 of No.~6. Thus, it remains to prove that for every closed, balanced, convex neighborhood \(\mathrm{V}\) of 0 in \(\mathrm{H}\) and for every equicontinuous subset \(\mathrm{N}\) of \(\mathcal{C}'(\mathrm{T})\) , there exists a neighborhood \(\mathrm{W}\) of 0 in \(\mathrm{G}\) such that the relations \(\mathbf{x} \in \mathrm{W}\) , \(\lambda \in \mathrm{N}\) imply that \(U(\lambda) \cdot \mathbf{x} \in \mathrm{V}\) . One can suppose that \(\mathrm{N} = \mathrm{S}^{\circ}\) , where \(\mathrm{S}\) is a neighborhood of 0 in \(\mathcal{C}(\mathrm{T})\) , consequently one can suppose that \(\mathrm{S}\) is the set of functions \(g \in \mathcal{C}(\mathrm{T})\) such that \(|g(t)| \leqslant 1\) on a compact subset \(\mathrm{K}\) of \(\mathrm{T}\) . It suffices to show that \(|\langle U(\lambda) \cdot \mathbf{x}, \mathbf{x}' \rangle| \leqslant 1\) for \(\mathbf{x} \in W\) , \(\mathbf{x}' \in V^\circ\) and \(\lambda \in S^\circ\) . Now, since \(U(K)\) is equicontinuous, there exists a neighborhood \(W\) of 0 in \(G\) such that the relations \(t \in K\) , \(\mathbf{x} \in W\) imply \(U(t) \cdot \mathbf{x} \in V\) ; the relations \(\mathbf{x} \in W\) , \(\mathbf{x}' \in V^\circ\) therefore imply that the function \(t \mapsto \langle U(t) \cdot \mathbf{x}, \mathbf{x}' \rangle\) belongs to \(S\) , hence that \(|\langle U(t) \cdot \mathbf{x}, \mathbf{x}' \rangle| = |\int\langle U(t) \cdot \mathbf{x}, \mathbf{x}' \rangle d\lambda(t)| \leqslant 1\) by the definition of \(S^\circ\) .

Let us now assume that \(U\) is a continuous mapping of \(\mathrm{T}\) into \(\mathrm{F}\) and, moreover, that \(U(\mathrm{T})\) is equicontinuous. Then, the same reasoning as above shows (since \(\mathrm{H}\) is quasi-complete) that for every bounded positive measure \(\mu\) on \(\mathrm{T}\) , \(\int U d\mu \in \mathrm{F}\) . One can therefore define, as above, a linear mapping \(\lambda \mapsto \int U d\lambda = U(\lambda)\) of \(\mathcal{M}^1(\mathrm{T})\) into \(\mathrm{F}\) that extends the analogous mapping of \(\mathcal{C}'(\mathrm{T})\) into \(\mathrm{F}\) . Moreover, for every closed, balanced, convex neighborhood \(\mathrm{V}\) of 0 in \(\mathrm{H}\) , there exists by hypothesis a neighborhood \(\mathrm{W}\) of 0 in \(\mathrm{G}\) such that for every \(\mathbf{x} \in \mathrm{W}\) and every \(t \in \mathrm{T}\) , one has \(U(t) \cdot \mathbf{x} \in \mathrm{V}\) , consequently (since \(\mathrm{V}\) is weakly closed) \(\int (U(t) \cdot \mathbf{x}) d\lambda(t) \in \| \lambda \| \cdot \mathrm{V}\) (No.~2, Prop. 5). In other words:

PROPOSITION 17. --- Let G, H be two Hausdorff locally convex spaces, with H assumed to be quasi-complete. Let U be a mapping of T into \(\mathcal{L}(G;H)\) such that \((t,\mathbf{x})\mapsto U(t)\cdot \mathbf{x}\) is continuous on T×G, and \(U(T)\) is equicontinuous. Then, if \(\mathcal{M}^1 (T)\) is equipped with its Banach space topology, the bilinear mapping \((\lambda ,\mathbf{x})\mapsto U(\lambda)\cdot \mathbf{x}\) of \(\mathcal{M}^1 (T)\times G\) into H is continuous (which implies, in particular, that the linear mapping \(\lambda \mapsto U(\lambda)\) of \(\mathcal{M}^1 (T)\) into \(\mathcal{L}(G;H)\) is continuous when \(\mathcal{L}(G;H)\) is equipped with the topology of bounded convergence).

PROPOSITION 18. --- Let \(G_{1}, G_{2}, H_{1}, H_{2}\) be four Hausdorff locally convex spaces, with \(H_{1}\) and \(H_{2}\) assumed to be quasi-complete. Let \(A : G_{1} \to G_{2}\) and \(B : H_{1} \to H_{2}\) be two continuous linear mappings. Let \(U_{1} : T \to \mathcal{L}(G_{1}; H_{1})\) , \(U_{2} : T \to \mathcal{L}(G_{2}; H_{2})\) be two mappings satisfying the conditions of Prop. 16 (resp. Prop. 17), and suppose that for every \(t \in T\) , \(B \circ U_{1}(t) = U_{2}(t) \circ A\) . Then, for every measure with compact support (resp. every bounded measure) \(\lambda\) on T, \(B \circ U_{1}(\lambda) = U_{2}(\lambda) \circ A\) .

Indeed, for every \(x \in G_{1}\) , one has (No.~1, Prop. 1)

\[
\begin{array}{c} \left(B \circ U _ {1} (\lambda)\right) \cdot \mathbf {x} = \int \Big (\left(B \circ U _ {1} (t)\right) \cdot \mathbf {x} \Big) d \lambda (t) \\ = \int \Big (\left(U _ {2} (t) \circ A\right) \cdot \mathbf {x} \Big) d \lambda (t) = U _ {2} (\lambda) \cdot (A \cdot \mathbf {x}). \end{array}
\]

Remarks. --- 1) Suppose that G and H are Banach spaces, and let U be a mapping of T into \(\mathcal{L}(\mathrm{G};\mathrm{H})\) such that \((t,\mathbf{x})\mapsto U(t)\cdot\mathbf{x}\) is continuous on \(T\times G\) . Note that this implies that the finite function \(t\mapsto\|U(t)\|\) is bounded on every compact subset of T and lower semi-continuous on T, being the upper envelope of the continuous functions \(t \mapsto |U(t) \cdot \mathbf{x}|\) as x runs over the ball \(|x| \leqslant 1\) in G. Set \(h(t) = \|U(t)\|\) . Then, for every positive measure \(\mu\) on T such that h is \(\mu\) -integrable, we again have \(\int U d\mu \in \mathcal{L}(\mathrm{G};\mathrm{H})\) . For, the measure \(\nu = h \cdot \mu\) is bounded by hypothesis; there therefore exists a partition of T formed by a \(\nu\) -negligible set N and a sequence \((\mathrm{K}_{n})\) of compact subsets. The argument made at the beginning of this No., applied to the measure \(\varphi_{K_{n}} \cdot \mu\) , shows that

\[
A _ {n} = \int \varphi_ {\mathrm{K} _ {n}} U d \mu \in \mathrm{F} = \mathcal {L} (\mathrm{G}; \mathrm{H}),
\]

and, moreover (No.~2, Prop. 6), \(\| A_n\| \leqslant \int \varphi_{\mathrm{K}_n}\| U\| d\mu \leqslant \nu (\mathrm{K}_n)\) . The series with general term \(A_{n}\) is therefore absolutely convergent in the Banach space \(\mathcal{L}(\mathrm{G};\mathrm{H})\) , and it is immediate that its sum is \(\int Ud\mu\) and that \(\left\| \int Ud\mu \right\| \leqslant \int \| U\| d\mu\) .

\begin{enumerate}
\def\labelenumi{\arabic{enumi})}
\setcounter{enumi}{1}
\tightlist
\item
  Suppose that \(\mathrm{G} = \mathrm{H}\) is quasi-complete and that \(U\) satisfies the hypotheses of Prop. 16. Let \(\mathbf{M}\) be a dense subset of the space \(\mathcal{C}'(\mathrm{T})\) , for the weak topology \(\sigma\big(\mathcal{C}'(\mathrm{T}),\mathcal{C}(\mathrm{T})\big)\) , and let \(\mathrm{X}\) be a closed linear subspace of \(\mathrm{H}\) such that \(U(\lambda)(\mathrm{X})\subset \mathrm{X}\) for every measure \(\lambda \in \mathrm{M}\) . Then also \(U(t)(\mathrm{X})\subset \mathrm{X}\) for every \(t\in \mathrm{T}\) : indeed, for every \(\mathbf{x}\in \mathrm{X}\) and every \(\mathbf{x}'\in \mathrm{H}'\) orthogonal to \(\mathrm{X}\) , by hypothesis \(\langle U(\lambda)\cdot \mathbf{x},\mathbf{x}'\rangle = 0\) for all \(\lambda \in \mathrm{M}\) , which may be written \(\int \langle U(t)\cdot \mathbf{x},\mathbf{x}'\rangle d\mu (t) = 0\) . The continuous function \(t\mapsto \langle U(t)\cdot \mathbf{x},\mathbf{x}'\rangle\) , being orthogonal to \(\mathrm{M}\) , is therefore 0, which yields \(\langle U(t)\cdot \mathbf{x},\mathbf{x}'\rangle = 0\) for every \(\mathbf{x}'\in \mathrm{X}^\circ\) , whence \(U(t)\cdot \mathbf{x}\in \mathrm{X}\) for all \(t\in \mathrm{T}\) and \(\mathbf{x}\in \mathrm{X}\) , and this proves our assertion.
\end{enumerate}

\section{VECTORIAL MEASURES}

\subsection{Definition of a vectorial measure}

The definition of a measure given in Ch. III, §1, No.~3 may be generalized as follows:

DEFINITION 1. --- Let F be a Hausdorff locally convex space over R. One calls vectorial measure on T with values in F every continuous linear mapping of the space \(\mathcal{H}(T)\) into F.

Def. 1 may also be expressed as follows: a vectorial measure on T with values in F is a linear mapping m of \(\mathcal{K}(T)\) into F such that, for every compact subset K of T, the restriction of m to \(\mathcal{K}(T,K)\) is continuous for the topology of uniform convergence. If \(f \in \mathcal{K}(T)\) , one also writes \(\int f dm\) or \(\int f(t) dm(t)\) instead of \(\mathbf{m}(f)\) . The measures with values in R are sometimes called real measures (Ch. III, §1, No.~5) or scalar measures on T.

Examples. --- 1) The identity mapping of \(\mathcal{H}(\mathrm{T})\) is a vectorial measure on \(\mathrm{T}\) with values in \(\mathcal{H}(\mathrm{T})\) .

\begin{enumerate}
\def\labelenumi{\arabic{enumi})}
\setcounter{enumi}{1}
\tightlist
\item
  *Let H be a complex Hilbert space, \(\mathcal{L}(\mathrm{H})\) the normed algebra of continuous endomorphisms of H. Let A be a subalgebra of \(\mathcal{L}(\mathrm{H})\) , commutative, closed, self-adjoint and containing the identity; one shows that there exist a compact space X and an isomorphism of the normed algebra A onto the algebra \(\mathcal{K}_{\mathbf{C}}(\mathrm{X}) = \mathcal{C}_{\mathbf{C}}(\mathrm{X})\) of continuous complex functions on X, equipped with the norm \(\| f \| = \sup_{x \in \mathrm{X}} |f(x)|\) . The inverse isomorphism, restricted to \(\mathcal{K}(\mathrm{X})\) , is a vectorial measure \(\mathbf{m}\) on X, with values in \(\mathcal{L}(\mathrm{H})\) , such that \(\mathbf{m}(fg) = \mathbf{m}(f)\mathbf{m}(g)\) .*
\end{enumerate}

Remarks. --- 1) For a linear mapping m of \(\mathcal{K}(\mathrm{T})\) into F to be a vectorial measure, it is necessary and sufficient that, for every compact subset K of T, the image under m of the unit ball \(\|f\|\leqslant1\) of \(\mathcal{K}(\mathrm{T},\mathrm{K})\) be bounded in F. The notion of vectorial measure with values in F is therefore the same for all the Hausdorff locally convex topologies on F that admit the same bounded sets, and in particular for all the topologies compatible with the duality between F and F' (TVS, IV, §1, No.~1, Prop. 1).

\begin{enumerate}
\def\labelenumi{\arabic{enumi})}
\setcounter{enumi}{1}
\item
  Let \(T_{1}\) be a locally compact space, \(F_{1}\) a Hausdorff locally convex space over R, u a continuous linear mapping of \(\mathcal{K}(T_{1})\) into \(\mathcal{K}(T)\) , and v a continuous linear mapping of F into \(F_{1}\) . If m is a vectorial measure on T with values in F, then \(v \circ m \circ u\) is a vectorial measure on \(T_{1}\) with values in \(F_{1}\) . In particular, if g is a continuous, finite numerical function on T, then \(f \mapsto \mathbf{m}(gf)\) is a vectorial measure on T with values in F, which is denoted \(g \cdot m\) ; if h is a second continuous, finite numerical function on T, then \(g \cdot (h \cdot \mathbf{m}) = (gh) \cdot \mathbf{m}\) .
\item
  Since the space \(\mathcal{H}(\mathrm{T})\) is the direct limit of the Banach spaces \(\mathcal{H}(\mathrm{T},\mathrm{K})\) , and is in particular barreled (TVS, III, §4, No.~1, Cor. of Prop. 2 and Cor. 3 of Prop. 3), in order that a linear mapping \(\mathbf{m}\) of \(\mathcal{H}(\mathrm{T})\) into \(\mathbf{F}\) be a vectorial measure, it is necessary and sufficient that, for every \(\mathbf{z}' \in \mathbf{F}'\) , \(\mathbf{z}' \circ \mathbf{m}\) be a scalar measure on \(\mathrm{T}\) (TVS, II, §6, No.~4, Prop. 5 and IV, §1, No.~3, Prop. 7).
\item
  In view of Remark 1, Prop. 1 of Ch. III, §2, No.~1 and its proof are again valid for vectorial measures. One can therefore define the support of a vectorial measure \(\mathbf{m}\) on \(T\) to be the complement of the largest open set \(U \subset T\) such that the restriction of \(\mathbf{m}\) to \(U\) is zero.
\end{enumerate}

\subsection{Integration with respect to a vectorial measure}

Let \(\mathbf{m}\) be a vectorial measure on \(T\) , with values in \(F\) . For every \(\mathbf{z}' \in F'\) , the mapping \(\mathbf{z}' \circ \mathbf{m}\) is a scalar measure on \(T\) , depending linearly on \(\mathbf{z}'\) . If \(f\) is a numerical function defined on \(T\) , we shall say, by an abuse of language, that the pair \((f, \mathbf{m})\) has the property \(\boldsymbol{P}\) if, for every \(\mathbf{z}' \in F'\) , the pair \((f, |\mathbf{z}' \circ \mathbf{m}|)\) has the property \(\boldsymbol{P}\) . For example, we shall say that \(f\) is essentially integrable for \(\mathbf{m}\) if, for every \(\mathbf{z}' \in F'\) , the function \(f\) is essentially integrable for \(|\mathbf{z}' \circ \mathbf{m}|\) . It comes to the same to say that \(f\) is essentially integrable for each of the measures \((\mathbf{z}' \circ \mathbf{m})^+\) and \((\mathbf{z}' \circ \mathbf{m})^-\) (Ch. V, §2, No.~2, Cor. 2 of Prop. 3).

PROPOSITION 1. --- Let m be a vectorial measure on T with values in F, f a numerical function on T that is essentially integrable for m. The mapping

\[
\mathbf {z} ^ {\prime} \mapsto \int f d (\mathbf {z} ^ {\prime} \circ \mathbf {m}) ^ {+} - \int f d (\mathbf {z} ^ {\prime} \circ \mathbf {m}) ^ {-}
\]

is a linear form on \(\mathbf{F}'\) .

Denoting this mapping by \(\Phi\) , it is immediate that \(\Phi(\lambda\mathbf{z}') = \lambda\Phi(\mathbf{z}')\) for all \(\lambda \in \mathbf{R}\) . Everything comes down to showing that \(\Phi(\mathbf{y}' + \mathbf{z}') = \Phi(\mathbf{y}') + \Phi(\mathbf{z}')\) . Set \(\mu = |\mathbf{y}' \circ \mathbf{m}| + |\mathbf{z}' \circ \mathbf{m}|\) ; by the Lebesgue-Nikodym theorem, one can then write \(\mathbf{y}' \circ \mathbf{m} = g \cdot \mu\) and \(\mathbf{z}' \circ \mathbf{m} = h \cdot \mu\) , where \(g\) and \(h\) are two bounded and locally \(\mu\) -integrable numerical functions (Ch. V, §5, No.~5, Th. 2); moreover, \((\mathbf{y}' \circ \mathbf{m})^+ = g^+ \cdot \mu\) and \((\mathbf{y}' \circ \mathbf{m})^- = g^- \cdot \mu\) , and the analogous relations hold with \(\mathbf{y}'\) replaced by \(\mathbf{z}'\) (resp. \(\mathbf{y}' + \mathbf{z}'\) ) and \(g\) by \(h\) (resp. \(g + h\) ). This being so, it is immediate that \(f\) is essentially \(\mu\) -integrable (Ch. V, §2, No.~2, Cor. 1 of Prop. 3), and the relation to be proved reduces to \((g + h)^+ - (g + h)^- = (g^+ - g^-) + (h^+ - h^-)\) , which is obvious.

DEFINITION 2. --- Let m be a vectorial measure on T with values in F, f a numerical function on T that is essentially integrable for m. One calls integral of f with respect to m, and denotes by \(\mathbf{m}(f)\) or \(\int f dm\) or again \(\int f(t) dm(t)\) , the element of \(F'^{*}\) defined by

\[
\left\langle \mathbf {z} ^ {\prime}, \int f d \mathbf {m} \right\rangle = \int f d (\mathbf {z} ^ {\prime} \circ \mathbf {m}) ^ {+} - \int f d (\mathbf {z} ^ {\prime} \circ \mathbf {m}) ^ {-}.\tag{1}
\]

We observe that if \(f \in \mathcal{K}(T)\) , the element \(\int f dm\) so defined coincides with the element denoted likewise in No.~1, because the second member of (1) is then \(\int f d(\mathbf{z}' \circ \mathbf{m}) = \mathbf{z}'(\mathbf{m}(f))\) by definition. Moreover, if in particular one applies Def. 2 to the case that F = R, one sees that for every \(z' \in F'\) , \(f\) is essentially integrable for the real measure \(\mathbf{z}' \circ \mathbf{m}\) , and that the second member of (1) may be written \(\int f d(\mathbf{z}' \circ \mathbf{m})\) .

Suppose now that \(f\) is essentially integrable for \(\mathbf{m}\) , and let \(\mathbf{z}' \in \mathbf{F}'\) . Set \(\mu = |\mathbf{z}' \circ \mathbf{m}|\) ; by the Lebesgue-Nikodym theorem, one can write \(\mathbf{z}' \circ \mathbf{m} = g \cdot \mu\) , where \(g\) is locally \(\mu\) -integrable and \(\| g \| \leqslant 1\) , and the proof of Prop. 1 shows that \(\int f d(\mathbf{z}' \circ \mathbf{m}) = \int fg d\mu\) . Consequently,

\[
\left| \int f d (\mathbf {z} ^ {\prime} \circ \mathbf {m}) \right| \leqslant \int | f | d | \mathbf {z} ^ {\prime} \circ \mathbf {m} |.\tag{2}
\]

It is clear that the set of finite numerical functions essentially integrable for m is a vector space over R; we shall denote by \(\mathcal{L}(\mathbf{m})\) this space equipped with the coarsest locally convex topology making continuous all of the linear forms \(f \mapsto \int f d(\mathbf{z}' \circ \mathbf{m})\) , where \(z'\) runs over \(F'\) . Note that in general the locally convex space \(\mathcal{L}(\mathbf{m})\) is not Hausdorff.

Example. --- Let us take for m the identity mapping of \(\mathcal{H}(\mathrm{T})\) onto itself. Since the dual of \(\mathcal{H}(\mathrm{T})\) is the space \(\mathcal{M}(\mathrm{T})\) of scalar measures on T, the functions \(f \in \mathcal{L}(\mathbf{m})\) are those that are essentially integrable for every scalar measure \(\mu\) (cf.~Exer. 1), and the integral \(\int f dm\) is the linear form \(\mu \mapsto \int f d\mu\) on \(\mathcal{M}(\mathrm{T})\) . One cannot have \(\int f d\mu = 0\) for every measure \(\mu \in \mathcal{M}(\mathrm{T})\) unless \(f = 0\) , as one sees on taking \(\mu = \varepsilon_t\) , where \(t\) is arbitrary in T; in other words, the mapping \(f \mapsto \int f dm\) is an injection of \(\mathcal{L}(\mathbf{m})\) into the algebraic dual of \(\mathcal{M}(\mathrm{T})\) , which extends the identity mapping of \(\mathcal{H}(\mathrm{T})\) . The relation \(\int f dm \in \mathrm{F} = \mathcal{H}(\mathrm{T})\) is therefore equivalent to \(f \in \mathcal{H}(\mathrm{T})\) .

Let \(u\) be a continuous linear mapping of \(F\) into a Hausdorff locally convex space \(G\) , and let us denote again by \(u\) its extension by bitransposition to a linear mapping of \(F'^*\) into \(G'^*\) (§1, No.~1). With this convention:

PROPOSITION 2. --- Every numerical function \(f\) essentially integrable for \(\mathbf{m}\) is essentially integrable for \(u \circ \mathbf{m}\) , and \(\int f d(u \circ \mathbf{m}) = u\left(\int f dm\right)\) .

The proposition is obvious, in view of the equality \(y'\circ u\circ m = {}^{t}u(y')\circ m\) for all \(y \in G'\) .

In general, if \(f \in \mathcal{L}(\mathbf{m})\) , the integral \(\int f dm\) belongs to \(\mathbf{F}'*\) but not to \(\mathbf{F}\) (see the above Example). However:

PROPOSITION 3. --- If the image under m of the set of \(f \in \mathcal{K}(\mathbf{T})\) such that \(\sup_{t \in \mathbf{T}} |f(t)| \leqslant 1\) is weakly relatively compact in F, then \(\int f dm \in \mathbf{F}\) for every bounded numerical function \(f\) essentially integrable for m.

Let \(\mathbf{A}\) be the set of \(f\in \mathcal{L}(\mathbf{m})\) such that \(\sup_{t\in \mathrm{T}}|f(t)|\leqslant 1\) , and let

\(B = A \cap \mathcal{K}(T)\) ; by hypothesis, \(\mathbf{m}(B)\) is weakly relatively compact in \(F\) , therefore it suffices to show that \(\mathbf{m}(A)\) is contained in the closure (in \(F'^*\) ) of \(\mathbf{m}(B)\) for the topology \(\sigma(F'^*, F')\) ; since \(\mathbf{m}(B)\) is convex and balanced, it suffices to prove that the polar of \(\mathbf{m}(\mathrm{B})\) in \(\mathrm{F}'\) is contained in that of \(\mathbf{m}(\mathrm{A})\) (TVS, II, §6, No.~3, Th. 1). Now, for a linear form \(\mathbf{z}' \in \mathrm{F}'\) to belong to \((\mathbf{m}(\mathrm{B}))^\circ\) , it is necessary and sufficient that \(|\langle \mathbf{z}', \mathbf{m}(g) \rangle| = |\int g d(\mathbf{z}' \circ \mathbf{m})| \leqslant 1\) for every function \(g \in \mathrm{B}\) , which signifies that the scalar measure \(|\mathbf{z}' \circ \mathbf{m}|\) is bounded and of norm \(\leqslant 1\) (Ch. III, §1, No.~8); but by (2) the latter condition implies that \(|\langle \mathbf{z}', \mathbf{m}(f) \rangle| \leqslant 1\) for every function \(f \in \mathrm{A}\) , whence \(\mathbf{z}' \in (\mathbf{m}(\mathrm{A}))^\circ\) .

COROLLARY 1. --- If, for every compact subset K of T, the image under m of the set of \(f \in \mathcal{K}(\mathrm{T},\mathrm{K})\) such that \(\sup_{t \in \mathrm{T}} |f(t)| \leqslant 1\) is weakly relatively compact in F, then \(\int f dm \in \mathrm{F}\) for every function \(f \in \mathcal{L}(\mathbf{m})\) that is bounded and has compact support, and \(\int f dm \in \mathrm{F}''\) for every function \(f \in \mathcal{L}(\mathbf{m})\) .

The first assertion may be deduced immediately from Prop. 3: if \(f\) is bounded and has compact support, and if \(U\) is a relatively compact open neighborhood of the support of \(f\) , then the restriction of \(\mathbf{m}\) to the subspace \(\mathcal{K}(U)\) is a measure \(\mathbf{m}_U\) on \(U\) that satisfies the conditions of Prop. 3, and \(\int f dm_U = \int f dm\) (Ch. V, §7, No.~1, Th. 1), therefore \(\int f dm \in F\) .

Now let \(f\) be any element of \(\mathcal{L}(\mathbf{m})\) ; for every compact subset \(K\) of \(T\) and every integer \(n > 0\) , let \(f_{n,K}\) be the numerical function on \(T\) defined as follows: if \(t \notin K\) , \(f_{n,K}(t) = 0\) ; if \(t \in K\) and \(|f(t) \leqslant n\) , \(f_{n,K}(t) = f(t)\) ; finally, if \(t \in K\) and \(|f(t)| > n\) , \(f_{n,K}(t) = nf(t)/|f(t)|\) . It is clear that for every \(t \in T\) , \(f(t)\) is the limit of \(f_{n,K}(t)\) with respect to the product filter of the Fréchet filter by the section filter of the (increasing directed) ordered set of compact subsets of \(T\) . Since \(|f_{n,K}| \leqslant |f|\) , it follows from Lebesgue's theorem and Prop. 10 of Ch. V, §1, No.~3, applied to each scalar measure \(|\mathbf{z}' \circ \mathbf{m}|\) , that \(f_{n,K}\) converges to \(f\) in \(\mathcal{L}(\mathbf{m})\) with respect to the preceding filter. Consequently, the integral \(\int f dm\) is in the closure in \(F'^*\) (for the topology \(\sigma(F'^*, F')\) ) of the set \(M\) of \(m(f_{n,K})\) . But the first part of the corollary shows that \(M \subset F\) , and, on the other hand, for every \(z' \in F'\) one has \(|\langle z', m(f_{n,K})\rangle| \leqslant \int |f| d|\mathbf{z}' \circ \mathbf{m}|\) , which shows that \(M\) is bounded in \(F_\sigma\) , hence also in \(F\) (TVS, IV, §1, No.~1, Prop. 1). Lemma 1 of §1, No.~2 therefore shows that \(\int f dm \in F''\) .

COROLLARY 2. --- If F is semi-reflexive, then \(\int f dm \in F\) for every numerical function \(f\) essentially integrable for \(\mathbf{m}\) .

\subsection{Majorizable vectorial measures}

Let \(q\) be a lower semi-continuous semi-norm on \(F\) . We shall denote by \(A_q'\) the set of \(z' \in F'\) such that \(|\langle z', x \rangle| \leqslant q(x)\) for all \(x \in F\) . This is the polar in \(\mathbf{F}'\) of the set of \(\mathbf{x} \in \mathbf{F}\) such that \(q(\mathbf{x}) \leqslant 1\) ; for every \(\mathbf{x} \in \mathbf{F}\) , \(q(\mathbf{x}) = \sup_{\mathbf{z}' \in \hat{\mathrm{A}}_q'} |\langle \mathbf{x}, \mathbf{z}' \rangle|\) .

DEFINITION 3. --- Let m be a vectorial measure on T with values in F. If q is a lower semi-continuous semi-norm on F, m is said to be q-majorizable if there exists a positive measure \(\mu\) such that \(|z^{\prime} \circ m| \leqslant \mu\) for every \(z^{\prime} \in A_{q}^{\prime}\) ; the supremum of the measures \(|z^{\prime} \circ m|\) as \(z^{\prime}\) runs over \(A_{q}^{\prime}\) (Ch. III, §2, No.~4, Th. 3) is then denoted \(q(\mathbf{m})\) . One says that m is majorizable if it is q-majorizable for every continuous semi-norm q on F.

If \(\mathbf{m}\) and \(\mathbf{m}'\) are both \(q\) -majorizable, it is immediate that \(\mathbf{m} + \mathbf{m}'\) is also \(q\) -majorizable and that

\[
q (\mathbf {m} + \mathbf {m} ^ {\prime}) \leqslant q (\mathbf {m}) + q (\mathbf {m} ^ {\prime}).
\]

When F is a normed space, with norm denoted \(|x|\) , to say that m is majorizable therefore means that the measures \(|z'\circ m|\) , where \(|z'| \leqslant 1\) , are bounded above \(^{1}\) by a same positive measure; one then denotes by \(|m|\) the supremum of this family of measures.

If \(\mathbf{F} = \mathbf{R}\) , the measure \(|m|\) corresponding to the Euclidean norm \(|x|\) on \(\mathbf{R}\) coincides with the measure denoted by \(|m|\) in Ch. III, §1, No.~5.

PROPOSITION 4. --- Let \((\mathrm{F}_i)_{1 \leqslant i \leqslant n}\) be a finite family of Hausdorff locally convex spaces, \(\mathbf{F} = \prod_{i=1}^{n} \mathbf{F}_i\) their product, \(q_i\) ( \(1 \leqslant i \leqslant n\) ) a lower semi-continuous semi-norm on \(\mathbf{F}_i\) , and \(q\) the semi-norm on \(\mathbf{F}\) defined by \(q(\mathbf{x}_i, \ldots, \mathbf{x}_n) = \sum_{i=1}^{n} q_i(\mathbf{x}_i)\) . If \(\mathbf{m}_i\) ( \(1 \leqslant i \leqslant n\) ) is a vectorial measure on \(\mathbf{T}\) with values in \(\mathbf{F}_i\) that is \(q_i\) -majorizable, then the measure \(\mathbf{m} = (\mathbf{m}_1, \ldots, \mathbf{m}_n)\) with values in \(\mathbf{F}\) is \(q\) -majorizable.

For, the dual \(\mathbf{F}'\) may be identified with \(\prod_{i=1}^{n} \mathrm{F}_i'\) , in such a way that if \(\mathbf{x} = (\mathbf{x}_i) \in \mathrm{F}\) , \(\mathbf{z}' = (\mathbf{z}_i') \in \mathrm{F}'\) , then \(\langle \mathbf{x}, \mathbf{z}' \rangle = \sum_{i=1}^{n} \langle \mathbf{x}_i, \mathbf{z}_i' \rangle\) . If \(|\langle \mathbf{x}, \mathbf{z}' \rangle| \leqslant q(\mathbf{x})\) for every \(\mathbf{x} \in \mathrm{F}\) , then in particular \(|\langle \mathbf{x}_i, \mathbf{z}_i' \rangle| \leqslant q_i(\mathbf{x}_i)\) for \(1 \leqslant i \leqslant n\) , and the converse is obvious, which shows that the set \(\mathrm{A}_q'\) is the product of the \(\mathrm{A}_{q_i}'\) . Since by hypothesis \(|\mathbf{z}_i' \circ \mathbf{m}_i| \leqslant q_i(\mathbf{m}_i)\) for \(\mathbf{z}_i' \in \mathrm{A}_{q_i}'\) , it follows that

\[
\left| \mathbf {z} ^ {\prime} \circ \mathbf {m} \right| \leqslant \sum_ {i = 1} ^ {n} \left| \mathbf {z} _ {i} ^ {\prime} \circ \mathbf {m} _ {i} \right| \leqslant \sum_ {i = 1} ^ {n} q _ {i} (\mathbf {m} _ {i})
\]

for every \(\mathbf{z}' \in \mathrm{A}_q'\) , which proves the proposition.

COROLLARY. --- If the space F is finite-dimensional, then every vectorial measure m with values in F is majorizable. In order that a numerical function be essentially integrable for m, it is necessary and sufficient that it be essentially integrable for \(|m|\) (where \(|x|\) denotes any norm on F).

PROPOSITION 5. --- Let q be a lower semi-continuous semi-norm on F. Let m be a q-majorizable measure, and let f be a function essentially integrable for m and such that \(\int f dm \in F\) . Then

\[
q \left(\int f d \mathbf {m}\right) \leqslant \int^ {\bullet} | f | d q (\mathbf {m}).
\]

For,

\[
q \left(\int f d \mathbf {m}\right) = \sup _ {\mathbf {z} ^ {\prime} \in \mathrm{A} _ {q} ^ {\prime}} \left| \left\langle \mathbf {z} ^ {\prime}, \int f d \mathbf {m} \right\rangle \right| \leqslant \sup _ {\mathbf {z} ^ {\prime} \in \mathrm{A} _ {q} ^ {\prime}} \int | f | d | \mathbf {z} ^ {\prime} \circ m | \leqslant \int^ {\bullet} | f | d q (\mathbf {m})
\]

by virtue of (1) and the relation \(|\mathbf{z}' \circ \mathbf{m}| \leqslant q(\mathbf{m})\) for \(\mathbf{z}' \in \mathrm{A}_q'\) .

PROPOSITION 6. --- Let F be a quasi-complete, Hausdorff locally convex space, m a majorizable measure on T with values in F. If f is a numerical function essentially integrable for all of the measures \(q(\mathbf{m})\) (where q runs over the set of continuous semi-norms on F), then f is essentially integrable for m, and \(\int f dm \in F\) .

We shall make use of the following auxiliary result. Let \((\mu_{\iota})_{\iota\in\mathrm{I}}\) be an increasing directed family of positive measures on T. Let us denote by \(\mathcal{L}^{1}\big((\mu_{\iota})_{\iota\in\mathrm{I}}\big)\) the vector space of finite numerical functions on T, essentially \(\mu_{\iota}\) -integrable for every \(\iota\in I\) , equipped with the topology defined by the semi-norms \(f\mapsto\mu_{\iota}(|f|)\;\ (\iota\in\mathrm{I})\) . Let \(L_{0}\) be the linear subspace of \(\mathcal{L}^{1}\big((\mu_{\iota})_{\iota\in\mathrm{I}}\big)\) generated by the products \(g\varphi_{K}\) , where g runs over the set of continuous finite numerical functions on T, and K over the set of compact subsets of T.

Lemma 1. --- When \(\mathcal{L}_0\) and \(\mathcal{K}(\mathrm{T})\) are equipped with the topology induced by that of \(\mathcal{L}^1\big((\mu_\iota)_{\iota \in \mathrm{I}}\big)\) :

\begin{enumerate}
\def\labelenumi{\alph{enumi})}
\item
  each element of \(\mathcal{L}_0\) is in the closure of some bounded subset of \(\mathcal{K}(\mathrm{T})\) ;
\item
  each element of \(\mathcal{L}^1\big((\mu_\iota)_{\iota \in \mathrm{I}}\big)\) is in the closure of some bounded subset of \(\mathcal{L}_0\) .
\end{enumerate}

To prove a), we may restrict ourselves to the case of an element of the form \(f = g\varphi_{\mathrm{K}}\) ( \(g \in \mathcal{C}(\mathrm{T})\) ), K compact in T). It is immediate (by virtue of Urysohn's theorem) that \(f\) is in the closure of the set B of functions of the form \(gh\) , where \(h\) describes the set of continuous mappings of T into {[}0, 1{]} that are equal to 1 on K and to 0 on the complement of a fixed compact neighborhood H of K. Moreover, the set B is bounded, because \(\mu_{\iota}(|gh|) \leqslant \mu_{\iota}(|f\varphi_{\mathrm{H}}|)\) for every function h having the preceding properties.

Let us now prove b); we may restrict ourselves to the case of an element \(f \geqslant 0\) of \(\mathcal{L}^1((\mu_\iota)_{\iota \in \mathrm{I}})\) . For every \(\iota \in \mathrm{I}\) and every \(\varepsilon > 0\) , there exists a compact subset \(\mathrm{K}(\iota, \varepsilon)\) of \(\mathrm{T}\) such that the restriction of \(f\) to \(\mathrm{K}(\iota, \varepsilon)\) is continuous and \(|\mu_\iota(|f - f\varphi_{\mathrm{K}(\iota, \varepsilon)}|) \leqslant \varepsilon\) . It is clear that \(f\) is in the closure of the set \(\mathrm{C}\) of \(f\varphi_{\mathrm{K}(\iota, \varepsilon)}\) (where \(\iota \in \mathrm{I}, \varepsilon > 0\) ). By virtue of Urysohn's theorem, the set \(\mathrm{C}\) in contained in \(\mathcal{L}_0\) ; moreover, it is bounded, because \(\mu_\varkappa(f\varphi_{\mathrm{K}(\iota, \varepsilon)}) \leqslant \mu_\varkappa(f)\) for all \(\iota \in \mathrm{I}, \varkappa \in \mathrm{I}\) and \(\varepsilon > 0\) .

Let us now prove Prop. 6: for every function \(g \in \mathcal{K}(\mathrm{T})\) and every continuous seminorm \(q\) on \(\mathrm{F}\) , \(q\big(\int g dm\big) \leqslant \int |g| d\big(q(\mathbf{m})\big)\) (Prop. 5), which implies that the mapping \(g \mapsto \int g dm\) of \(\mathcal{K}(\mathrm{T})\) into \(\mathrm{F}\) is continuous when \(\mathcal{K}(\mathrm{T})\) is equipped with the topology induced by that of \(\mathcal{L}^1\big((q(\mathbf{m}))_{q \in \mathbb{Q}}\big)\) (Q the set of continuous semi-norms on \(\mathrm{F}\) ). Consequently, by the preceding lemma and Prop. 10 of TVS, III, §1, No.~6, this mapping may be extended by continuity, first to a continuous linear mapping \(v_0\) of \(\mathcal{L}_0\) into \(\mathrm{F}\) , then to a continuous linear mapping \(v\) of \(\mathcal{L}^1\big((q(\mathbf{m}))_{q \in \mathbb{Q}}\big)\) into \(\mathrm{F}\) . Moreover, for every \(\mathbf{z}' \in \mathrm{F}'\) the relation \(\langle \mathbf{z}', v(f) \rangle = \int f d(\mathbf{z}' \circ \mathbf{m})\) holds, by the definition of \(v\) , for every \(f \in \mathcal{K}(\mathrm{T})\) ; since \(|\mathbf{z}' \circ \mathbf{m}| \leqslant q(\mathbf{m})\) for \(q(\mathbf{z}) = |\langle \mathbf{z}', \mathbf{z} \rangle|\) , the mapping \(f \mapsto \int f d(\mathbf{z}' \circ \mathbf{m})\) is continuous on \(\mathcal{L}^1\big((q(\mathbf{m}))_{q \in \mathbb{Q}}\big)\) , therefore again by continuity, the relation \(\langle \mathbf{z}', v(f) = \int f d(\mathbf{z}' \circ \mathbf{m})\) holds for every function \(f \in \mathcal{L}^1\big((q(\mathbf{m}))_{q \in \mathbb{Q}}\big)\) . It follows that \(\nu(f) = \int f dm\) , which completes the proof.

\subsection{\texorpdfstring{Vectorial measures with base \(\mu\)}{Vectorial measures with base mu}}

DEFINITION 4. --- Let \(\mu\) be a positive measure on T. A vectorial measure m on T, with values in F, is said to be a measure with base \(\mu\) if there exists a mapping f of T into F, scalarly locally \(\mu\) -integrable, such that \(\mathbf{m}(g) = \int g\mathbf{f}d\mu\) for every function \(g \in \mathcal{K}(\mathrm{T})\) . One then says that f is a density of m with respect to \(\mu\) , and one writes \(\mathbf{m} = \mathbf{f} \cdot \mu\) .

It is immediate that if \(f_{1}\) and \(f_{2}\) are two densities of m with respect to \(\mu\) , then \(f_{1}-f_{2}\) is scalarly locally \(\mu\) -negligible (Ch. V, §5, No.~3, Cor. 2 of Prop. 3); recall that in general this does not imply that \(f_{1}-f_{2}\) is zero locally almost everywhere (cf.~§1, Exer. 12 and No.~1, Remark 2).

Let \(\mathbf{m}\) be a measure with base \(\mu\) , of density \(\mathbf{f}\) . For a numerical function \(g\) to be essentially integrable for \(\mathbf{m}\) , it is necessary and sufficient that \(gf\) be scalarly essentially \(\mu\) -integrable (Ch. V, §5, No.~3, Th. 1).

PROPOSITION 7. --- Let f be a mapping scalarly locally integrable with respect to a positive measure \(\mu\) on T, such that, for every function \(g \in \mathcal{K}(T)\) , one has \(\int g\mathbf{f} d\mu \in F\) . Then the mapping \(g \mapsto \int g\mathbf{f} d\mu\) of \(\mathcal{K}(T)\) into F is a vectorial measure on T, with base \(\mu\) and density f with respect to \(\mu\) .

For (No.~1, Remark 3), it suffices to show that, setting \(\mathbf{m}(g) = \int g\mathbf{f}d\mu\) , \(\mathbf{z}' \circ \mathbf{m}\) is a scalar measure for every \(\mathbf{z}' \in \mathrm{F}'\) . But since \(\mathbf{z}'(\mathbf{m}(g)) = \int g\langle \mathbf{z}', \mathbf{f} \rangle d\mu\) , one has \(\mathbf{z}' \circ \mathbf{m} = \langle \mathbf{z}', \mathbf{f} \rangle \cdot \mu\) , whence our assertion.

PROPOSITION 8. --- Let \(\mu\) be a positive measure on T, m a measure on T with values in F, with base \(\mu\) and density f with respect to \(\mu\) . Let q be a lower semi-continuous semi-norm on F.

\begin{enumerate}
\def\labelenumi{\alph{enumi})}
\item
  If, for every compact subset \(\mathrm{K}\) of \(\mathrm{T}\) , the upper integral \(\int_{\mathrm{K}}^{*}(q \circ \mathbf{f}) d\mu\) is finite, then \(\mathbf{m}\) is \(q\) -majorizable.
\item
  If \(\mathbf{m}\) is \(q\) -majorizable, then \(q(\mathbf{m})\) has base \(\mu\) ; if in addition \(\mathbf{f}\) is \(\mu\) -measurable as a mapping of \(\mathrm{T}\) into \(\mathrm{F}_{\sigma}\) , then \(q \circ \mathbf{f}\) is locally \(\mu\) -integrable and \(q(\mathbf{m}) = (q \circ \mathbf{f}) \cdot \mu\) .
\end{enumerate}

\begin{enumerate}
\def\labelenumi{\alph{enumi})}
\tightlist
\item
  For every finite subset \(J\) of \(A_q'\) , let us denote by \(\lambda_J\) the supremum of the measures \(|\mathbf{z}' \circ \mathbf{m}|\) , where \(\mathbf{z}'\) runs over \(J\) ; if \(g_J = \sup_{\mathbf{z}' \in J} |\langle \mathbf{z}', \mathbf{f} \rangle|\) then \(\lambda_{J}=g_{J}\cdot\mu\) (Ch. V, §5, No.~2, Prop. 2). For every relatively compact open subset U of T, let \(\lambda_{J,U}\) be the restriction of \(\lambda_{J}\) to U; let us first show that as J runs over the directed set F of finite subsets of \(A_{q}^{\prime}\) , the family \((\lambda_{J,U})\) is bounded above in \(\mathcal{M}(U)\) . Indeed, for every function \(h\geqslant0\) in \(K(U)\) ,
\end{enumerate}

\[
\int h d \lambda_ {\mathrm{J}, \mathrm{U}} = \int h g _ {\mathrm{J}} d \mu \leqslant \int^ {*} (q \circ \mathbf {f}) h d \mu \leqslant \| h \| \int_ {\mathrm{U}} ^ {*} (q \circ \mathbf {f}) d \mu ,
\]

whence our assertion (Ch. II, §2, No.~2). Let \(\nu_{\mathrm{U}}\) be the supremum of this family of measures in \(\mathcal{M}(\mathrm{U})\) . If \(\mathrm{U}'\) is a second relatively compact open subset of T such that \(\mathrm{U} \subset \mathrm{U}'\) , then \(\nu_{\mathrm{U}}\) is the restriction of \(\nu_{\mathrm{U}'}\) to U, as follows immediately from the expression of the supremum of an increasing directed set of measures (Ch. II, §2, No.~2) and the fact that \(\lambda_{\mathrm{J,U}}\) is the restriction to U of \(\lambda_{\mathrm{J,U}'}\) . Thus there is one and only one positive measure \(\nu\) whose restriction to each U is \(\nu_{\mathrm{U}}\) (Ch. III, §2, No.~1, Prop. 1), and it is clear that \(\nu = q(\mathbf{m})\) .

\begin{enumerate}
\def\labelenumi{\alph{enumi})}
\setcounter{enumi}{1}
\tightlist
\item
  Since the measures \(\lambda_{\mathrm{J}}\) have base \(\mu\) , so does their supremum \(q(\mathbf{m})\) (Ch. V, §5, No.~5, Th. 2). If \(\mathbf{f}\) is \(\mu\) -measurable for the topology \(\sigma(\mathrm{F},\mathrm{F}')\) on \(\mathrm{F}\) , it follows at once from the definitions that the mapping \(g:t\mapsto(g_{\mathrm{J}}(t))_{\mathrm{J}\in\mathfrak{F}}\) of \(\mathrm{T}\) into the product space \(\mathbf{R}^{\mathfrak{F}}\) is \(\mu\) -measurable. The restriction of \(q\circ\mathbf{f}=\sup_{\mathrm{J}\in\mathfrak{F}}g_{\mathrm{J}}\) to every compact subset of \(\mathrm{T}\) on which \(g\) is continuous is lower semi-continuous; consequently \(q \circ f\) is \(\mu\) -measurable (Ch. IV, §5, No.~5, Cor. of Prop. 8 and No.~10, Prop. 15). Let \(K\) be a compact subset of T; it admits a partition consisting of a \(\mu\) -negligible set and a sequence \((\mathrm{K}_n)\) of compact sets on which \(g\) is continuous. Then \(\int_{\mathrm{K}_n}^*\left(q\circ \mathbf{f}\right)d\mu = \sup_{\mathrm{J}}\int_{\mathrm{K}_n}g_{\mathrm{J}}d\mu \leqslant \int_{\mathrm{K}_n}dq(\mathbf{m})\) for all \(n\) (Ch. IV, §1, No.~1, Th. 1 and Ch. V, §7, No.~1, Prop. 1), whence \(\int_{\mathrm{K}}^{*}(q\circ \mathbf{f})d\mu = \sum_{n}\int_{\mathrm{K}_n}^*\left(q\circ \mathbf{f}\right)d\mu \leqslant \int_{\mathrm{K}}dq(\mathbf{m})\) . But this proves that \(q\circ \mathbf{f}\) is locally \(\mu\) -integrable and that \(\lambda_{\mathrm{J}}\leqslant (q\circ \mathbf{f})\cdot \mu \leqslant q(\mathbf{m})\) for all \(\mathbf{J}\in \mathfrak{F}\) ; whence, by definition, \(q(\mathbf{m}) = (q\circ \mathbf{f})\cdot \mu\) .
\end{enumerate}

Remark. --- Suppose that there exists in \(A_{q}^{\prime}\) a countable subset D that is dense for \(\sigma(\mathbf{F}^{\prime},\mathbf{F})\) ; then, the function \(q\circ f\) is always \(\mu\) -measurable, because \(q(\mathbf{f}(t)) = \sup_{\mathbf{z}^{\prime}\in\mathrm{D}} |\langle\mathbf{z}^{\prime},\mathbf{f}(t)\rangle|\) (Ch. IV, §5, No.~4, Cor. 1 of Th. 2). Then, for every

compact subset K of T, \(\int_{K}^{*}(q\circ\mathbf{f})d\mu=\sup_{J}\int_{K}g_{J}d\mu\) , where J runs over the countable directed set of finite subsets of D (Ch. IV, §1, No.~1, Cor. of Th. 3); thus, one sees that in this case the condition that \(\int_{K}^{*}(q\circ\mathbf{f})d\mu<+\infty\) for every compact subset K of T is necessary and sufficient for m to be q-majorizable.

PROPOSITION 9. --- Let F be a finite-dimensional Banach space. Every measure m on T with values in F is a measure with base \textbar m\textbar. If \(m = f \cdot |m|\) , then \(|f(t)| = 1\) locally almost everywhere for \textbar m\textbar. For m to have base \(\mu\) , it is necessary and sufficient that \textbar m\textbar{} have base \(\mu\) , and if \(m = g \cdot \mu\) then \(|m| = |g| \cdot \mu\) .

Let \((\mathbf{e}_i)_{1 \leqslant i \leqslant n}\) and \((\mathbf{e}_i')_{1 \leqslant i \leqslant n}\) be dual bases of F and \(\mathbf{F}'\) (A, II, §2, No.~6) with \(|\mathbf{e}_i'| = 1\) for all \(i\) . Then \(|\mathbf{e}_i' \circ \mathbf{m}| \leqslant |\mathbf{m}|\) for every index \(i\) , therefore (Ch. V, §5, No.~5, Th. 2) \(\mathbf{e}_i' \circ \mathbf{m} = h_i \cdot |\mathbf{m}|\) , where \(h_i\) is bounded and \(|\mathbf{m}|\) -measurable. Setting \(\mathbf{h} = \sum_{i=1}^{n} h_i \cdot \mathbf{e}_i\) , we therefore have \(\mathbf{m} = \mathbf{h} \cdot |\mathbf{m}|\) . If \(\mathbf{m} = \mathbf{f} \cdot |\mathbf{m}|\) , Prop. 8 shows that \(|\mathbf{m}| = |\mathbf{f}| \cdot |\mathbf{m}|\) , whence \(|\mathbf{f}(t)| = 1\) locally almost everywhere for \(|\mathbf{m}|\) (Ch. V, §5, No.~3, Cor. 2 of Prop. 3). The final assertion follows at once from Prop. 8.

Remark. --- If \(z = \sum_{i=1}^{n} z_i e_i\) then \(\psi(z_1, \ldots, z_n) = |z|\) is a positively homogeneous continuous function on \(R to the n\) . Setting \(\mu_i = e_i' \circ m = h_i \cdot |m|\) , by definition (Ch. V, §5, No.~9) \(\psi(\mu_1, \ldots, \mu_n) = \psi(h_1, \ldots, h_n) \cdot |m| = |h| \cdot |m| = |m|\) .

\subsection{The Dunford--Pettis theorem}

Let \(\mu\) be a positive measure on T. A vectorial measure \(\mathbf{m}\) on T, with values in F, is said to be scalarly of base \(\mu\) (or to have base \(\mu\) scalarly) if, for every \(\mathbf{z}' \in \mathrm{F}'\) , the scalar measure \(\mathbf{z}' \circ \mathbf{m}\) has base \(\mu\) . If a vectorial measure m with values in F has base \(\mu\) , then it has base \(\mu\) scalarly: for, if \(m = f \cdot \mu\) then \(z' \circ m = \langle z', f \rangle \cdot \mu\) for every \(z' \in F'\) . But there exist vectorial measures that are scalarly of base \(\mu\) without having base \(\mu\) (Exer. 17), and, on the other hand, there exist vectorial measures that are not scalarly of base \(\nu\) for any positive measure \(\nu\) ; note however that every majorizable measure m with values in a normed space is scalarly of base \(|m|\) , by virtue of the Lebesgue--Nikodym theorem.

Example. --- Let us take for m the identity mapping of \(\mathcal{K}(\mathrm{T})\) . To say that m is scalarly of base \(\mu\) means that every real measure on T has base \(\mu\) . In particular, the point measure \(\varepsilon_{t}\) ( \(t \in T\) ) must have base \(\mu\) , which requires that \(\mu(\{t\}) > 0\) for every \(t \in T\) , and implies in particular that every compact subset of T is countable.

In this No.~we are going to prove a result that generalizes one of the consequences of the Lebesgue--Nikodym theorem, namely, that the dual of \(\mathrm{L}^{1}(\mu)\) is \(\mathrm{L}^{\infty}(\mu)\) (Ch. V, §5, No.~8, Th. 4), and that gives a sufficient condition for a vectorial measure that is scalarly of base \(\mu\) to have base \(\mu\) .

Let \(\pi\) be the canonical mapping of \(\mathcal{L}^{\infty}(\mu)\) onto \(\mathrm{L}^{\infty}(\mu)\) . We shall say that a linear subspace \(\mathbf{G}\) of \(\mathrm{L}^{\infty}\) has the lifting property if there exists a linear mapping \(\rho\) of \(\mathbf{G}\) into \(\mathcal{L}^{\infty}(\mu)\) (called a lifting of \(\mathbf{G}\) ) such that \(\pi \circ \rho\) is the identity on \(\mathbf{G}\) and \(|\rho(f)(t)| \leqslant \mathrm{N}_{\infty}(f)\) for all \(t \in \mathrm{T}\) and \(f \in \mathrm{G}\) .

One proves that if \(\mu\) is Lebesgue measure on \(\mathbf{R}^n\) , the entire space \(\mathrm{L}^{\infty}(\mathbf{R}^{n},\mu)\) has the lifting property (Exer. 18).

Lemma 2. --- Every separable subspace G of the Banach space \(\mathrm{L}^{\infty}(\mathrm{T},\mu)\) has the lifting property.

By the hypothesis, there exists a countable dense subset H of G that is a linear subspace with respect to the field Q of rational numbers; let \((h_{n})\) be a (countable) basis of H over Q. For every integer n, let \(h_{n}^{\prime}\) be an element of \(L^{\infty}\) such that \(\pi(h_{n}^{\prime}) = h_{n}\) , and let \(\rho^{\prime}\) be the Q-linear mapping of H into \(L^{\infty}\) defined by \(\rho^{\prime}(h_{n}) = h_{n}^{\prime}\) ; it is clear that \(\pi \circ \rho^{\prime}\) is the identity on H. Moreover, for every \(h \in H\) one has \(|\rho^{\prime}(h)(t)| \leqslant N_{\infty}(h)\) except at the points t of a locally negligible set A(h). Let A be the union of the A(h) for \(h \in H\) , which is also locally negligible. For every \(h \in H\) , denote by \(\rho(h)\) the function \(h^{\prime\prime} \in L^{\infty}\) such that \(h^{\prime\prime}(t) = \rho^{\prime}(h)(t)\) if \(t \notin A\) , and \(h^{\prime\prime}(t) = 0\) if \(t \in A\) . It is clear that \(\rho\) is a Q-linear mapping of H into the subspace B of bounded functions in \(L^{\infty}\) , such that \(\pi \circ \rho\) is the identity on H and such that \(|\rho(h)(t)| \leqslant N_{\infty}(h)\) for all \(h \in H\) and \(t \in T\) . Since B is a Banach space for the norm \(\|f\| = \sup_{t \in T} |f(t)|\) (Ch. IV, §5, No.~4, Th. 2),

\(\rho\) may be extended to a continuous R-linear mapping, again denoted \(\rho\) , of G into B, which is obviously a lifting of G.

DEFINITION 5. --- Let F be a Hausdorff locally convex space, \(F_{s}^{\prime}\) its dual equipped with the topology \(\sigma(F^{\prime},F)\) . We denote by \(L_{F_{s}^{\prime}}^{\infty}\) the vector space of mappings f of T into \(F_{s}^{\prime}\) , such that f is scalarly \(\mu\) -measurable and is equal scalarly locally almost everywhere (for \(\mu\) ) to a mapping of T into an equicontinuous subset of \(F^{\prime}\) . We denote by \(L_{F_{s}^{\prime}}^{\infty}\) the quotient space of \(L_{F_{s}^{\prime}}^{\infty}\) by the space of scalarly locally \(\mu\) -negligible mappings of T into \(F_{s}^{\prime}\) .

When F satisfies the hypotheses of §1, No.~5, Prop. 13, the functions in \(\mathcal{L}_{\mathrm{F}_s^\prime}^\infty\) are \(\mu\) -measurable for the weak topology \(\sigma(\mathrm{F}',\mathrm{F})\) , but are not necessarily measurable for the strong topology on \(\mathrm{F}'\) , even if F is a Banach space (§1, Exer. 17). Under the same conditions, the scalarly locally \(\mu\) -negligible mappings of T into \(\mathrm{F}_s^\prime\) are identical to the locally \(\mu\) -negligible mappings of T into \(\mathrm{F}_s^\prime\) (§1, No.~1, Remark 2).

When F is a separable normed space, the elements of \(L_{F_{s}^{\prime}}^{\infty}\) are the mappings f of T into \(F_{s}^{\prime}\) such that f is scalarly \(\mu\) -measurable and \(|f|\) is bounded in measure; one can then define a normed space structure on the space \(L_{F_{s}^{\prime}}^{\infty}\) , by equipping it with the norm \(N_{\infty}\) (Ch. IV, §6, No.~3).

Lemma 3. --- Let F be a Hausdorff locally convex space, f an element of \(\mathcal{L}_{\mathrm{F}_s^\prime}^\infty\) . For every \(\mathbf{z} \in \mathrm{F}\) , one has \(\langle \mathbf{z}, \mathbf{f} \rangle \in \mathcal{L}^\infty\) , and the linear mapping \(z \mapsto \pi(\langle \mathbf{z}, \mathbf{f} \rangle)\) of F into \(\mathrm{L}^\infty\) is continuous; if, moreover, F is a normed space, then \(\mathrm{N}_\infty(\langle \mathbf{z}, \mathbf{f} \rangle) \leqslant |\mathbf{z}| \cdot \sup_{t \in \mathrm{T}} |\mathbf{f}(t)|\) .

It is clear by definition that \(\langle z,f\rangle\) is \(\mu\) -measurable and bounded in measure; replacing if necessary f by a function belonging to the same class of \(L_{F_{s}}^{\infty}\) , we can suppose in addition that \(\mathbf{f}(T)\subset V^{\circ}\) , where V is a balanced convex neighborhood of 0 in F (which does not modify \(\langle z,f\rangle\) except on a locally negligible set, depending on z). Then the relation \(z\in V\) implies that \(|\langle z,f(t)\rangle|\leqslant1\) for all \(t\in T\) , which proves the continuity of \(\mathbf{z}\mapsto\pi(\langle\mathbf{z},\mathbf{f}\rangle)\) . The final assertion is obvious.

Lemma 4. --- Let F be a Hausdorff locally convex space, f an element of \(\mathcal{L}_{\mathrm{F}_s}^\infty\) . For every numerical function \(g \in \overline{\mathcal{L}}^1\) , the function gf is scalarly essentially \(\mu\) -integrable and \(\int g\mathbf{f} d\mu \in \mathrm{F}'\) .

For every \(z \in F\) , \(\langle z, f \rangle\) belongs to \(L^{\infty}\) , whence the first assertion. One can suppose moreover, without modifying \(\int g f d\mu\) , that \(\mathbf{f}(T) \subset V^{\circ}\) , where V is a balanced convex neighborhood of 0 in F. Then the relation \(z \in V\) implies \(|\langle z, f(t) \rangle| \leqslant 1\) for all \(t \in T\) , whence \(|\langle z, \int g f d\mu \rangle| = |\int \langle z, f \rangle g d\mu| \leqslant \overline{N}_{1}(g)\) , which proves that \(\int g f d\mu \in F'\) .

THEOREM 1. --- Let F be a Hausdorff locally convex space that contains a countable dense subset. For every function \(\mathbf{f} \in \mathcal{L}_{\mathrm{F}_s}^\infty\) and every \(\mathbf{z} \in \mathbf{F}\) , let \(v_{\mathbf{f}}(\mathbf{z}) = \pi (\langle \mathbf{z}, \mathbf{f} \rangle) \in \mathrm{L}^{\infty}\) ; the mapping \(\mathbf{f} \mapsto v_{\mathbf{f}}\) defines, by passage to the quotient, a linear bijection of \(L_{F^{\prime}s}^{\infty}\) onto the space \(\mathcal{L}(F;L^{\infty})\) of continuous linear mappings of F into \(L^{\infty}\) . If F is a normed space, this bijection is an isometry.

In view of Lemma 3, the first assertion will be demonstrated if one proves that for every continuous mapping \(u\) of \(F\) into \(L^{\infty}\) , there exists a function \(\mathbf{f} \in \mathcal{L}_{F_s'}^{\infty}\) such that \(\pi(\langle \mathbf{z}, \mathbf{f} \rangle) = u(\mathbf{z})\) for all \(\mathbf{z} \in F\) , and that the class of \(\mathbf{f}\) in \(L_{F_s'}^{\infty}\) is uniquely determined by this condition. The second point is immediate, because if \(\pi(\langle \mathbf{z}, \mathbf{f} \rangle) = \pi(\langle \mathbf{z}, \mathbf{f}_1 \rangle)\) for all \(\mathbf{z} \in F\) , then \(\mathbf{f}_1 - \mathbf{f}\) is scalarly locally negligible. On the other hand, there exists a lifting \(\rho\) of \(u(F)\) into \(\mathcal{L}^{\infty}\) (Lemma 2). For every \(t \in T\) , the mapping \(\mathbf{z} \mapsto \rho(u(\mathbf{z}))(t)\) is a continuous linear form on \(F\) , that is, an element \(f(t)\) of \(F'\) . The function \(\mathbf{f}\) is scalarly \(\mu\) -measurable since \(\langle \mathbf{z}, \mathbf{f} \rangle = \rho(u(\mathbf{z})) \in \mathcal{L}^{\infty}\) for every \(\mathbf{z} \in F\) ; one has \(\pi(\langle \mathbf{z}, \mathbf{f} \rangle) = u(\mathbf{z})\) ; finally, for every \(t \in T\) and every \(\mathbf{z}\) belonging to the inverse image \(V\) under \(u\) of the unit ball of \(L^{\infty}\) ,

\[
\left| \langle \mathbf {z}, \mathbf {f} (t) \rangle \right| = \left| \rho (u (\mathbf {z})) (t) \right| \leqslant N _ {\infty} (u (\mathbf {z})) \leqslant 1,
\]

which shows that \(\mathbf{f}(t) \in \mathrm{V}^{\circ}\) for all \(t \in \mathrm{T}\) .

If, moreover, F is a normed space, the foregoing shows that

\[
\sup _ {t \in \mathrm{T}} | \mathbf {f} (t) | \leqslant \| u \|.
\]

But on the other hand (Lemma 3), \(\mathrm{N}_{\infty}(u(\mathbf{z})) \leqslant |\mathbf{z}| \cdot \sup_{t \in \mathrm{T}} |\mathbf{f}(t)|\) , and this inequality continues to hold when f is modified arbitrarily on a locally negligible set. It follows that \(\|u\| = \mathrm{N}_{\infty}(|\mathbf{f}|)\) .

COROLLARY 1. --- Let F be a Hausdorff locally convex space containing a countable dense subset. For every function \(\mathbf{f} \in \mathcal{L}_{\mathrm{F}_s}^\infty\) , every \(\mathbf{z} \in \mathrm{F}\) and every function \(g \in \mathcal{L}^1\) , let \(\Phi_{\mathbf{f}}(\mathbf{z}, \widetilde{g}) = \int \langle \mathbf{z}, \mathbf{f}(t) \rangle g(t) d\mu(t)\) . The mapping \(\mathbf{f} \mapsto \Phi_{\mathbf{f}}\) defines, by passage to the quotient, a linear bijection of \(\mathrm{L}_{\mathrm{F}_s}^\infty\) onto the space \(\mathcal{B}(\mathrm{F}, \mathrm{L}^1)\) of continuous bilinear forms on \(\mathrm{F} \times \mathrm{L}^1\) . If \(\mathrm{F}\) is a normed space, this bijection is an isometry.

One can suppose that \(\mathbf{f}(\mathrm{T})\) is an equicontinuous subset of \(F'\) . It is then clear that \(\Phi_{f}\) is separately continuous, and, with the notations of Th. 1 and the Appendix, one has (taking into account the fact that \(L^{\infty}\) is the dual of \(L^{1}\) (Ch. V, §5, No.~8, Th. 4)) \(^{l}\Phi_{f}=v_{f}\) . The corollary then follows from Th. 1 above and from the Appendix, No.~1, Prop. 1 and its corollary.

COROLLARY 2 (Dunford--Pettis theorem). --- Let F be a Hausdorff locally convex space containing a countable dense subset. For every function \(f \in L_{F_{s}^{\infty}}\) and every function \(g \in L^{1}\) , let \(w_{\mathbf{f}}(\widetilde{g}) = \int g\mathbf{f} \, d\mu \in \mathbf{F}'\) (Lemma 4). The mapping \(f \mapsto w_{f}\) defines, by passage to the quotient, a linear bijection of \(L_{F_{s}^{\prime}}^{\infty}\) onto the space \(\mathcal{R}(L^{1}, F^{\prime})\) of linear mappings u of \(L^{1}\) into \(F^{\prime}\) such that the image under u of the unit ball of \(L^{1}\) is an equicontinuous subset of \(F^{\prime}\) . If F is a normed space (in which case \(\mathcal{R}(L^{1}, F^{\prime})\) is the space of continuous linear mappings of \(L^{1}\) into the strong dual of F), the bijection \(f \mapsto w_{f}\) is an isometry.

Taking into account the fact that \(L^{\infty}\) is the dual of \(L^{1}\) , this follows from the preceding corollary and from the Appendix, No.~1, Prop. 1 and its corollary.

Remark. --- It is clear that the mappings \(u \in \mathcal{R}(L^{1}, F')\) are continuous for every S-topology on \(F'\) (S a covering of F by bounded subsets). Conversely, if F is assumed moreover to be barreled, then every continuous linear mapping of \(L^{1}\) into \(F'\) equipped with an S-topology transforms the unit ball of \(L^{1}\) into a bounded subset of \(F'\) , which is therefore equicontinuous (TVS, III, §4, No.~2, Th. 1).

COROLLARY 3. --- Let F be a Hausdorff locally convex space containing a countable dense subset, m a vectorial measure on T with values in the weak dual F' of F. If the image under m of the set B of functions g in \(\mathcal{K}(T)\) such that \(\mu(|g|) \leqslant 1\) is contained in a closed and convex equicontinuous subset H' of F', then m has base \(\mu\) and there exists a density f of m with respect to \(\mu\) such that f(t) ∈ H' for all t ∈ T.

The hypothesis implies that m is continuous when \(\mathcal{K}(T)\) is equipped with the topology induced by that of \(\mathcal{L}^{1}(\mu)\) (defined by the semi-norm \(N_{1}\) ); it may therefore be extended to a continuous linear mapping u of \(\mathcal{L}^{1}(\mu)\) into the completion G of the weak dual of F; but since \(H'\) is a compact subset of G and the image under u of the set \(\overline{B}\) of \(f \in L^{1}\) such that \(\mathrm{N}_{1}(f) \leqslant 1\) is contained in the closure of \(H'\) in G, one has \(u(\overline{\mathrm{B}}) \subset \mathrm{H}'\) , therefore u maps \(L^{1}\) into \(F'\) . Since the relation \(\mathrm{N}_{1}(f) \leqslant \varepsilon\) implies that \(u(f) \in \varepsilon\mathrm{H}'\) , one has \(u(g) = 0\) if g is \(\mu\) -negligible, and Cor. 2 may therefore be applied to the mapping of \(L^{1}\) into \(F'\) obtained from u by passing to the quotient; whence the corollary.

COROLLARY 4. --- Let F be a separable normed space, and m a measure on T with values in the strong dual F', majorizable for the norm of F'. Then m is a measure with base \textbar m\textbar, and if \(m = f \cdot |m|\) then \(|f(t)| = 1\) locally almost everywhere for \textbar m\textbar.

By hypothesis, for every \(\mathbf{z} \in \mathbf{F}\) such that \(|\mathbf{z}| \leqslant 1\) , one has \(|\langle \mathbf{z}, \mathbf{m}(g) \rangle| \leqslant |\mathbf{m}|(|g|)\) for every function \(g \in \mathcal{K}(\mathrm{T})\) , consequently \(|\mathbf{m}(g)| \leqslant |\mathbf{m}|(|g|)\) (TVS, IV, §2, No.~4, formula (1)). Since every ball in \(\mathbf{F}'\) is equicontinuous, Cor. 3 is applicable and shows that \(\mathbf{m}\) has base \(|\mathbf{m}|\) ; moreover, if \(\mathbf{m} = \mathbf{f} \cdot |\mathbf{m}|\) then \(\mathbf{f}\) is \(|\mathbf{m}|\) -measurable for \(\sigma(\mathrm{F}', \mathrm{F})\) (§1, No.~5, Prop. 13) and \(|\mathbf{m}| = |\mathbf{f}| \cdot |\mathbf{m}|\) (No.~4, Prop. 8), which proves the corollary (Ch. V, §5, No.~3, Cor. 2 of Prop. 3).

If this corollary is applied to the case that \(\mathbf{F}\) if finite-dimensional, one recovers as a special case the first part of Prop. 9.

\subsection{\texorpdfstring{Dual of the space \(L_{F}^{1}\) (F a separable Banach space)}{Dual of the space L sub F superscript 1 (F a separable Banach space)}}

PROPOSITION 10. --- Let F be a separable Banach space. For every function \(\mathbf{f} \in \overline{\mathcal{L}}_{\mathrm{F}}^{1}\) and every function \(\mathbf{g} \in \mathcal{L}_{\mathrm{F}'_s}^{\infty}\) , the numerical function \(\langle \mathbf{f}, \mathbf{g} \rangle : t \mapsto \langle \mathbf{f}(t), \mathbf{g}(t) \rangle\) is essentially \(\mu\) -integrable, and

\[
\left| \int \langle \mathbf {f}, \mathbf {g} \rangle d \mu \right| \leqslant \overline {{{{N}}}} _ {1} (\mathbf {f}) N _ {\infty} (\mathbf {g}).\tag{3}
\]

For every class \(\dot{\mathbf{g}}\in\mathrm{L}_{\mathrm{F}^{\prime}s}^{\infty}\) , let \(\theta(\dot{\mathbf{g}})\) be the continuous linear form on \(\mathrm{L}_{\mathrm{F}}^{1}\) deduced from the linear form \(\mathbf{f}\mapsto\int\langle\mathbf{f},\mathbf{g}\rangle d\mu\) on \(\overline{\mathcal{L}}_{\mathrm{F}}^{1}\) by passage to the quotient; then \(\theta\) is a linear isometry of \(\mathrm{L}_{\mathrm{F}^{\prime}s}^{\infty}\) onto the strong dual \((\mathrm{L}_{\mathrm{F}}^{1})^{\prime}\) of the Banach space \(\mathrm{L}_{\mathrm{F}}^{1}\) .

For every compact subset K of T and every \(\varepsilon > 0\) , there exists a compact subset \(K'\) of K such that \(\mu(K - K') \leqslant \varepsilon\) and the restriction of f (resp. g) to \(K'\) is a continuous mapping of \(K'\) into F (resp. into \(F'_s\) ); consequently \(\mathbf{g}(K')\) is weakly compact, hence equicontinuous on \(F'\) (TVS, III, §4, No.~2, Th. 1 or IV, §1, Exer. 10). Now, the restriction of the canonical bilinear form on \(F \times F'\) to the product of F and an equicontinuous subset of \(F'\) is continuous for the product topology of the topology of F and \(\sigma(F', F)\) (GT, X, §2, No.~1, Cor. 4 of Prop. 1); it follows that the restriction of \(\langle f, g \rangle\) to \(K'\) is continuous, hence that \(\langle f, g \rangle\) is \(\mu\) -measurable. Moreover,

\[
| \langle \mathbf {f} (t), \mathbf {g} (t) \rangle | \leqslant | \mathbf {f} (t) | \cdot | \mathbf {g} (t) | \leqslant | \mathbf {f} (t) | N _ {\infty} (\mathbf {g})
\]

locally almost everywhere, consequently \(\langle f,g\rangle\) is essentially \(\mu\) -integrable and the inequality (3) holds (Ch. IV, §5, No.~6, Th. 5 and Ch. V, §1, No.~3, Lemma).

It remains to show that \(\theta\) is a surjective isometry. Let u be a continuous linear form on \(L_{F}^{1}\) . The mapping \((h,\mathbf{z})\mapsto u(\widetilde{h}\mathbf{z})\) is a continuous bilinear form on \(L^{1}\times F\) , because

\[
\left| u \left(\widetilde {h} \mathbf {z}\right) \right| \leqslant \| u \| \cdot \mathrm{N} _ {1} (h \mathbf {z}) \leqslant \| u \| \cdot | \mathbf {z} | \cdot \mathrm{N} _ {1} (h).
\]

By Cor. 1 of Th. 1 of No.~5, there exists a mapping g of T into F', belonging to \(L_{F_{s}^{\infty}}\) , such that \(|\mathbf{g}(t)| \leqslant \|u\|\) for all \(t \in T\) and such that \(u(\widetilde{h}\mathbf{z}) = \int \langle h\mathbf{z}, \mathbf{g} \rangle d\mu\) for every function \(h \in L^{1}\) with class \(\widetilde{h}\) in \(L^{1}\) and every \(z \in F\) . In other words, the linear forms u and \(\theta(\dot{\mathbf{g}})\) coincide on the subspace of \(L_{F}^{1}\) generated by the elements of the form \(\widetilde{h}\mathbf{z}\) ( \(\widetilde{h} \in L^{1}\) , \(z \in F\) ). Since this subspace is dense in \(L_{F}^{1}\) (Ch. IV, §3, No.~5, Prop. 10), it follows that \(u = \theta(\dot{\mathbf{g}})\) , which already proves that \(\theta\) is surjective. Moreover, by (3), \(\|\theta(\dot{\mathbf{g}})\| \leqslant N_{\infty}(\mathbf{g}) \leqslant \|u\| = \|\theta(\dot{\mathbf{g}})\|\) , whence \(\|\theta(\dot{\mathbf{g}})\| = N_{\infty}(\mathbf{g})\) , and this concludes the proof.
